% MobileCAD 中文使用手册 —— 用 XeLaTeX 编译：
%   xelatex main.tex && xelatex main.tex
\documentclass[UTF8,fontset=none,11pt,openany,oneside]{ctexbook}
% macOS 自带中文字体：宋体正文、黑体标题；fontset=none 避免 ctex 去找未安装的 STXihei
\setCJKmainfont{Songti SC}
\setCJKsansfont{Heiti SC}
\setCJKmonofont{Songti SC}
\setCJKfamilyfont{zhhei}{Heiti SC}
\setCJKfamilyfont{zhsong}{Songti SC}
\renewcommand{\CJKfamilydefault}{\CJKrmdefault}

\usepackage[a4paper,top=2.4cm,bottom=2.4cm,left=2.5cm,right=2.5cm,headheight=22pt]{geometry}
\usepackage{graphicx}
\usepackage{xcolor}
\usepackage{booktabs}
\usepackage{tabularx}
\usepackage{longtable}
\usepackage{array}
\usepackage{makecell}
\usepackage{enumitem}
\usepackage{listings}
\usepackage[most]{tcolorbox}
\usepackage{fancyhdr}
\usepackage{multicol}
\usepackage{caption}
\usepackage{fontawesome5}
\usepackage{amssymb}
\usepackage{seqsplit}
\usepackage{fontawesome5}
\usepackage{newunicodechar}
\newfontfamily\symfont{Apple Symbols}
% 宋体缺的符号：菜单箭头、勾叉、手势图标等，统一在此映射
\newunicodechar{✓}{\faCheck}
\newunicodechar{✗}{\faTimes}
\newunicodechar{✕}{\faTimes}
\newunicodechar{✎}{\faPen}
\newunicodechar{☰}{{\symfont ☰}}
\newunicodechar{▸}{{\symfont ▸}}
\newunicodechar{▾}{{\symfont ▾}}
\newunicodechar{◀}{{\symfont ◀}}
\newunicodechar{▶}{{\symfont ▶}}
\newunicodechar{▲}{{\symfont ▲}}
\newunicodechar{▼}{{\symfont ▼}}
\newunicodechar{↶}{{\symfont ↶}}
\newunicodechar{↷}{{\symfont ↷}}
\newunicodechar{↔}{{\symfont ↔}}
\newunicodechar{□}{\ensuremath{\square}}
\newunicodechar{≥}{\ensuremath{\geq}}
\newunicodechar{≤}{\ensuremath{\leq}}
\newunicodechar{ν}{\ensuremath{\nu}}
\newunicodechar{φ}{\ensuremath{\varphi}}
\usepackage[hidelinks,bookmarksnumbered=true]{hyperref}
\hypersetup{pdftitle={MobileCAD 中文使用手册},pdfauthor={MobileCAD 中文手册},pdfsubject={移动端参数化 CAD 教程},pdfkeywords={MobileCAD,FreeCAD,Android,CAD,教程,中文}}
% 中文段落里长英文/路径/命令 ID 的断行余量
\setlength{\emergencystretch}{2.5em}
\tolerance=3000

\graphicspath{{figures/}}

% ---------------------------------------------------------------- colors
\definecolor{mcblue}{HTML}{1651B0}
\definecolor{mclight}{HTML}{E8F0FB}
\definecolor{mcgray}{HTML}{5B6770}
\definecolor{mcorange}{HTML}{C25E00}
\definecolor{mcgreen}{HTML}{1E7B34}
\definecolor{mcred}{HTML}{A32020}
\definecolor{codebg}{HTML}{F6F7F9}

% ---------------------------------------------------------------- page style
\pagestyle{fancy}
\fancyhf{}
\fancyhead[L]{\small\color{mcgray}\leftmark}
\fancyhead[R]{\small\color{mcgray}MobileCAD 中文手册}
\fancyfoot[C]{\small\thepage}
\renewcommand{\chaptermark}[1]{\markboth{#1}{}}
\renewcommand{\headrulewidth}{0.4pt}
\renewcommand{\headrule}{\hbox to\headwidth{\color{mclight}\leaders\hrule height \headrulewidth\hfill}}

% ---------------------------------------------------------------- helpers
\newcommand{\ui}[1]{\textbf{#1}}                      % 界面上的按钮/标签
\newcommand{\cmd}[1]{\texttt{\small\seqsplit{#1}}}    % FreeCAD 命令 ID（长 ID 可断行）
\newcommand{\file}[1]{\texttt{\small\seqsplit{#1}}}   % 文件名/路径
\newcommand{\kbd}[1]{\colorbox{codebg}{\texttt{\small #1}}}
\newcommand{\menupath}[1]{\textsf{#1}}                % 菜单路径
\newcommand{\term}[1]{\textbf{\color{mcblue}#1}}      % 术语
\newcommand{\figref}[1]{图~\ref{#1}}

\captionsetup{font=small,labelfont={bf,color=mcblue},skip=6pt}
\setlist{itemsep=2pt,topsep=3pt,parsep=0pt}
\setcounter{tocdepth}{2}
\setcounter{secnumdepth}{2}

% ---------------------------------------------------------------- boxes
\newtcolorbox{tip}[1][提示]{
  enhanced,breakable,colback=mclight,colframe=mcblue!35,
  fonttitle=\bfseries\small,coltitle=white,colbacktitle=mcblue,
  title={#1},left=6pt,right=6pt,top=4pt,bottom=4pt,arc=2mm
}
\newtcolorbox{warn}[1][注意]{
  enhanced,breakable,colback=orange!6,colframe=mcorange!45,
  fonttitle=\bfseries\small,coltitle=white,colbacktitle=mcorange,
  title={#1},left=6pt,right=6pt,top=4pt,bottom=4pt,arc=2mm
}
\newtcolorbox{practice}[1][动手练习]{
  enhanced,breakable,colback=mcgreen!5,colframe=mcgreen!45,
  fonttitle=\bfseries\small,coltitle=white,colbacktitle=mcgreen,
  title={#1},left=6pt,right=6pt,top=4pt,bottom=4pt,arc=2mm
}

% ---------------------------------------------------------------- code
\lstdefinestyle{mc}{
  backgroundcolor=\color{codebg},basicstyle=\ttfamily\small,
  keywordstyle=\color{mcblue},stringstyle=\color{mcgreen},
  commentstyle=\color{mcgray}\itshape,numbers=none,
  frame=single,rulecolor=\color{mclight},framesep=3pt,
  breaklines=true,showstringspaces=false,columns=fullflexible,
  xleftmargin=2pt,xrightmargin=2pt,aboveskip=4pt,belowskip=4pt
}
\lstset{style=mc}

% ---------------------------------------------------------------- figures
\newcommand{\screenshot}[3][0.92]{%
  \begin{figure}[htbp]\centering
    \includegraphics[width=#1\linewidth]{#2}
    \caption{#3}
  \end{figure}}

\begin{document}

% ================================================================ title
\begin{titlepage}
\centering
\vspace*{2.2cm}
{\Huge\bfseries\color{mcblue} MobileCAD\par}
\vspace{0.5cm}
{\LARGE 中文使用手册\par}
\vspace{0.3cm}
{\large 从零开始，完全掌握移动端参数化 CAD\par}
\vspace{1.6cm}
\includegraphics[width=0.22\linewidth]{logo.jpg}\par
\vspace{1.6cm}
{\large 基于 FreeCAD 引擎 · Android 平板与手机\par}
\vspace{0.3cm}
{\large 适用于 MobileCAD 0.5.x（FreeCAD main e95a31c）\par}
\vfill
{\small 本手册与 MobileCAD 源码同步编写，示例均可在应用内完成。\par}
{\small 文档版本 1.0 · 2026 年 11 月\par}
\vspace{0.6cm}
{\small 许可证：GNU General Public License v3.0 or later\par}
\end{titlepage}

\frontmatter
\tableofcontents
\chapter*{About this manual (English)}
\addcontentsline{toc}{chapter}{About this manual (English)}
\markboth{About this manual}{About this manual}

This is a complete, task-oriented handbook for \textbf{MobileCAD}, the Android front end of the FreeCAD engine.
It is written from the source code — the Kotlin/Compose interface, the Python bridge and the bundled FreeCAD
modules — and takes a reader who has never opened the app to a finished parametric part, an assembly, a
drawing, a finite-element run, toolpaths and a BIM model.

\begin{itemize}
  \item \textbf{Scope:} 17 chapters and 3 appendices, about 90 pages. Chapters 1--4 get a beginner productive
        (UI tour, gestures, first part); 5--7 cover the modelling core (sketch, PartDesign, Part); 8--14 cover
        the remaining workbenches (Draft, TechDraw, expressions, Assembly, FEM, CAM, BIM, mesh, points);
        15--17 cover files, efficiency and troubleshooting; the appendices hold a command reference, a bilingual
        glossary and an FAQ.
  \item \textbf{Accuracy:} commands, dialog rows and status messages are taken from the source (the command
        catalog is generated from FreeCAD's own sources), and limitations are stated where the port has them
        (e.g. no Netgen mesher, no CAM material-removal simulation, no DWG import).
  \item \textbf{Language:} the manual is written in Simplified Chinese because it accompanies the Chinese
        interface; every UI label is given as it appears in Chinese (the English command name is shown in
        parentheses on first use), and Appendix B maps both.
  \item \textbf{Screenshots:} from a Samsung Galaxy Tab S9+ (Android 16) running the Chinese build, plus the
        project's own marketing screenshots.
  \item \textbf{Building:} \file{cd manual \&\& xelatex main.tex} (run twice for the table of contents). It needs
        XeLaTeX with \file{ctex}, and a CJK font — the preamble uses macOS' Songti SC / Heiti SC; on other
        systems change the four \file{\textbackslash setCJK...font} lines in \file{main.tex}.
  \item \textbf{License:} GNU General Public License v3.0 or later, like the application. Contributions and
        corrections are welcome as issues or pull requests against \file{manual/}.
\end{itemize}

\chapter*{写在前面}
\addcontentsline{toc}{chapter}{写在前面}
\markboth{写在前面}{写在前面}

MobileCAD 把 FreeCAD 的建模引擎搬到了手机上：同一套几何内核（Open CASCADE）、同一套参数化建模
体系、同一种文档格式（\file{.FCStd}），却配了一套为手指设计的界面。它不是“看图纸”的浏览器，而是
能真正建模、装配、出图、算有限元、生成刀路的完整 CAD。

这本手册的目标只有一个：\textbf{让你从第一次打开应用，到能独立完成一个中等复杂度的产品设计}。因此它
不是功能清单，而是按“先跑通、再深入、最后融会贯通”的顺序组织：

\begin{itemize}
  \item \textbf{第 1--4 章}：认识它、装上它、看懂界面，并完成第一个零件。读到这里你已经能画东西了。
  \item \textbf{第 5--7 章}：草图、零件设计、零件工作台——参数化建模的核心，占全书篇幅最大。
  \item \textbf{第 8--14 章}：其他工作台：绘图（Draft）、工程图（TechDraw）、变量与表达式、
        装配、有限元、CAM 与 BIM。
  \item \textbf{第 15--17 章}：文件交换、效率技巧、故障排除。
  \item \textbf{附录}：常用命令速查、中英术语对照、常见问题。
\end{itemize}

\section*{怎么读}
\begin{itemize}
  \item 第一次接触 CAD：按顺序读，每章的\emph{动手练习}都做一遍。
  \item 用过桌面 FreeCAD：第 3 章看交互差异，第 5--7 章看移动端特有的对话框与 Tip 机制，其余按需查。
  \item 只想解决问题：直接看目录、附录 A 的命令速查和附录 C 的常见问题。
\end{itemize}

\section*{约定}
\begin{itemize}
  \item \ui{粗体}表示界面上的按钮、标签或菜单项，例如 \ui{凸台}、\ui{更多选项}。
  \item \cmd{等宽字体}表示 FreeCAD 命令 ID，例如 \cmd{PartDesign\_Pad}，便于在源码与文档间对照。
  \item \menupath{菜单 → 子菜单}表示操作路径，例如 \menupath{文件 → 保存}。
  \item 本手册以\textbf{中文界面}为准，并在括号里给出英文原文，例如 \ui{凸台}（Pad）。切换语言的方法见 3.7 节。
\end{itemize}

\section*{关于这本手册}
本手册依据 MobileCAD 源码（应用层 Kotlin + Python 桥接 + FreeCAD 引擎）逐项核对写成，命令名、对话框
字段、提示语均取自源码。截图来自真实设备（三星 Galaxy Tab S9+，Android 16）与应用官方素材。若你发现
手册与当前版本不符，欢迎在项目仓库提 issue——CAD 软件演进很快，手册也会随之更新。


\mainmatter
\chapter{什么是 MobileCAD}

\section{一句话概括}

MobileCAD 是 \textbf{FreeCAD 引擎的 Android 前端}：它把桌面 FreeCAD 的建模内核、工作台与 Python 生态
打包进一个 APK，再用 Kotlin + Jetpack Compose 重写了整层用户界面——为触屏、为平板、也为手机。

这意味着两件事：

\begin{enumerate}
  \item 你得到的不是“简化版”建模器。参数化草图、实体建模、曲面、装配、工程图、有限元、CAM、BIM、
        网格与点云，背后都是 FreeCAD 原生的 C++ 实现；
  \item 文档格式与 FreeCAD 完全互通。手机上限时画的零件，可以在电脑上用 FreeCAD 打开继续改；反之亦然。
\end{enumerate}

\begin{tip}[一个引擎，两种界面]
桌面 FreeCAD 用 Qt 画界面；MobileCAD 用 Compose 画界面，但\emph{不}启动 FreeCAD 的 Qt 界面——它直接调用
FreeCAD 的 App 层（\file{FreeCAD.so}、\file{Part.so}、\file{PartDesign.so}……），并通过一个以 Python 实现的
“Qt 模拟层”让 FreeCAD 自带的任务面板（Pad 参数、BIM 墙参数、FEM 约束……）仍能运行，只是被渲染成了
Compose 的控件。所以你看到的每个对话框，逻辑都来自 FreeCAD 本身。
\end{tip}

\section{能力地图}

MobileCAD 复刻了 FreeCAD 的 \textbf{22 个工作台、984 个命令中的 980 个}（源码中的可比矩阵见
\file{research/PARITY.md}），未移植的只有桌面窗口管理一类没有手机对应物的命令。按用途归类：

\begin{center}
\small
\begin{tabularx}{\linewidth}{@{}lX@{}}
\toprule
\textbf{类别} & \textbf{工作台与能力} \\
\midrule
参数化建模 & \textbf{草图}（Sketcher）：几何、约束、尺寸、构造线、B 样条；\textbf{零件设计}（PartDesign）：
凸台、凹槽、旋转、放样、管道、螺旋、孔、圆角、倒角、拔模、厚度、阵列、镜像、参考几何 \\
实体与曲面 & \textbf{零件}（Part）：基本体、布尔、拉伸/旋转/扫掠/放样、偏移、检查、测量；\textbf{曲面}（Surface）\\
装配 & \textbf{装配}（Assembly）：13 种关节、接地、求解、冲突诊断、BOM、快照、仿真 \\
工程图 & \textbf{工程图}（TechDraw）：页面与模板、投影组、剖视、详图、尺寸、符号、导出 PDF/SVG/DXF \\
数据与表达式 & \textbf{变量}（Spreadsheet）：电子表格、别名、跨对象表达式、单位制 \\
仿真与制造 & \textbf{有限元}（FEM，CalculiX 求解器在设备上运行）；\textbf{CAM}（作业、刀具、刀路、后处理）\\
建筑与交换 & \textbf{BIM}（墙、板、柱、窗、楼梯、屋顶、IFC 导入导出）；\textbf{网格}、\textbf{点云}、
\textbf{OpenSCAD}、\textbf{逆向工程}、\textbf{机器人}、\textbf{检查}、\textbf{材质} \\
核心工具 & \textbf{Std}：文件、编辑、视图、测量、宏、依赖图、工程图导出、首选项等 \\
\bottomrule
\end{tabularx}
\end{center}

\begin{tip}[“全部命令”是否都常用？]
不是。984 个命令里有大量是同义或专业工作台的细分操作。日常建模真正高频的约 \textbf{40 个}：草图里的
十来个绘制/约束工具、零件设计里的凸台/凹槽/圆角/阵列、文件的打开保存导出。本书会把它们讲透，其余在
第 8 章与附录 A 中按图索骥。
\end{tip}

\section{与桌面 FreeCAD 的关系}

\begin{center}
\small
\begin{tabularx}{\linewidth}{@{}lXX@{}}
\toprule
 & \textbf{桌面 FreeCAD} & \textbf{MobileCAD} \\
\midrule
几何内核 & Open CASCADE & 相同（同一份交叉编译的库）\\
文档格式 & \file{.FCStd} & 相同，可直接互开 \\
Python 控制台 & 有 & 无（但宏与 Python 特征仍可运行）\\
任务面板 & Qt 对话框 & 同一套 Python 面板，渲染为触屏控件\\
界面 & Qt Widgets，双栏大量面板 & Compose，单工作区 + 浮动面板 \\
输入 & 鼠标键盘、快捷键 & 手指手势 + 数字键盘 + 软键盘；支持硬件键盘\\
绘图 / CAM / FEM & 完整 & 完整（求解器随 APK 打包）\\
桌面窗口管理 & 多窗口、停靠 & 无（移动端不需要）\\
\bottomrule
\end{tabularx}
\end{center}

\begin{warn}[版本对齐]
MobileCAD 的每个版本绑定一个 FreeCAD 提交。本手册对应 \textbf{FreeCAD main e95a31c}（应用内顶栏副标题
会显示，例如“零件 · 基于 FreeCAD main e95a31c”）。用桌面 FreeCAD 打开同一文档时，建议使用相近版本，
以免新特性（如某些 PartDesign 参数）在旧版里丢失。
\end{warn}

\section{运行条件与性能预期}

\begin{itemize}
  \item \textbf{系统}：Android 8.0 及以上，\textbf{arm64} 设备（手机、平板；Android 上的 Quest 头显走独立的
        MobileCAD XR 版本）。
  \item \textbf{存储}：安装包约 165 MB，首次启动解压运行时后占用约 1 GB；建议预留 2 GB 以上。
  \item \textbf{内存}：建模流畅建议 6 GB 以上；8--12 GB 平板可轻松处理几百个特征的装配。
  \item \textbf{性能}：实体建模与桌面 FreeCAD 同量级；有限元求解受手机散热限制，网格规模建议在桌面
        的 1/3 以内；大型装配的显示做了移动端优化（视口按需重绘、离屏缓存）。
\end{itemize}

\section{授权、隐私与数据}

MobileCAD 是 \textbf{GPL-3.0-or-later} 的自由软件，包含 FreeCAD（LGPL）、Open CASCADE（LGPL）、Qt（LGPL）、
Gmsh、CalculiX、IfcOpenShell 等组件。所有计算都在设备本地完成：\textbf{不上传文档、不依赖云}。唯一的
网络访问是你主动触发的“检查更新”（从 GitHub Releases 读取版本信息）以及 BIM/帮助里打开的网页。
完整许可列表在 \menupath{帮助 → 关于 MobileCAD → 开源许可}。

\chapter{安装与第一次启动}

\section{安装}

\subsection{官方渠道}

\begin{enumerate}
  \item \textbf{GitHub Releases}（推荐）：打开 \url{https://github.com/LucasCoraca/mobilecad/releases/latest}，
        下载 \file{MobileCAD-arm64.apk}，用系统文件管理器点开安装。首次安装需要允许该来源“安装未知应用”。
  \item \textbf{F-Droid}：项目提供 F-Droid 构建配方（\file{fdroid/}），在 F-Droid 里搜索 MobileCAD 即可。
        F-Droid 版本通过 F-Droid 自身更新，不使用应用内更新。
  \item \textbf{中文测试版}（本书配套）：由 \file{i18n/zh-hans} 分支构建，包名 \file{app.mobilecad.i18n}，
        可与官方版\textbf{并存}安装，界面为简体中文。它使用独立的测试签名，不参与官方更新。
\end{enumerate}

\begin{warn}[签名与升级]
官方 APK 与任何第三方重打包版本签名不同，\textbf{不能互相覆盖安装}。若你先装了官方版，再想换其他构建，
需要先卸载（文档保存在存储卡上不受影响；应用内的设置与最近列表会丢失）。
\end{warn}

\subsection{首次启动会发生什么}

第一次启动时，应用要做三件事，通常几十秒内完成（平板约 20--40 秒）：

\begin{enumerate}
  \item \textbf{解压运行时}：把 APK 内约 1 GB 的 FreeCAD、Python、Open CASCADE、Qt 运行库解压到应用私有目录；
  \item \textbf{启动引擎}：加载 100 多个本地库、初始化 Python 3.14、导入 FreeCAD，启动画面会显示
        \ui{正在启动 CAD 引擎…}（英文界面为 \ui{Starting the CAD engine…}）；应用升级后的第一次启动还会先显示
        \ui{正在解压 CAD 引擎…}（Unpacking the CAD engine…）；
  \item \textbf{建立空文档}：进入一个名为 \ui{未命名}（Unnamed）的新文档，特征树为空。
\end{enumerate}

之后每次启动只需 2--5 秒。解压只发生一次（应用升级后若运行时版本变化会重新解压）。若引擎启动失败，
启动画面会以红字显示 \ui{CAD 引擎启动失败：…}，并保留错误原因便于反馈。

\begin{tip}[不需要存储权限]
MobileCAD \textbf{不申请存储权限}。打开、保存、导入、导出全部通过 Android 的系统文件选择器（SAF）完成：
你选择哪个目录，应用就能读写哪个文件。文档默认保存在你选的位置，应用私有目录只存配置与崩溃恢复副本。
\end{tip}

\section{三十秒上手}

启动后你会看到如下的界面（中文版）：

\screenshot{ui-main-zh.jpg}{MobileCAD 中文主界面。左侧是特征树，顶部是 ☰ 菜单与文档信息，中间是 3D 视图，底部是元素选项卡（零件/装配/工程图/变量/更多），视图右上角是导航立方体。}

\begin{enumerate}
  \item \textbf{转一转}：单指在视图上拖动——模型绕中心旋转。
  \item \textbf{放大}：双指捏合缩放；双指拖动平移。
  \item \textbf{换个方向}：点右上角导航立方体的 \ui{前}/\ui{上}/\ui{右} 面，视图会动画转到对应方向；
        拖动立方体还能自由旋转。
  \item \textbf{打开菜单}：点左上角 ☰，能看到 文件、编辑、视图、工具、宏、帮助 六组菜单。
  \item \textbf{画点东西}：点顶部工具条的 \ui{草图}，在展开的工具栏里选 \ui{新建草图}，挑一个基准面，
        即可进入草图模式（第 4 章会完整走一遍）。
\end{enumerate}

\section{手势速查}

MobileCAD 的手势在“建模”和“草图”两种状态下有所不同，这一点和桌面 CAD 差别最大，值得先记住：

\begin{center}
\small
\begin{tabularx}{\linewidth}{@{}lXX@{}}
\toprule
\textbf{手势} & \textbf{建模状态} & \textbf{草图状态} \\
\midrule
单指拖动 & 旋转视图 & 绘制（手指上方 50 dp 处为实际落点）\\
双指拖动 & 平移 & 平移（第二指在 300 ms 内落下视为取消绘制）\\
双指捏合 & 缩放 & 缩放 \\
三指拖动 & 以三指中心旋转（大范围翻转） & 同左 \\
单指轻触 & 选择面/边/顶点 & 选择几何/约束 \\
双击 & 清除选择 & 清除选择 \\
双指轻触 & 打开右键菜单（等同桌面右键） & 同左 \\
长按 & 精确指针：准星出现在手指上方 60 dp，便于不遮挡目标 & 长按尺寸标签可编辑数值 \\
双指长按 & 框选（拖动矩形选择对象） & — \\
\bottomrule
\end{tabularx}
\end{center}

\begin{tip}[为什么指针在手指上方？]
手指会遮住落点，这是移动端 CAD 最大的交互难题。MobileCAD 在绘制与长按精确模式下，把实际落点画在手指
上方一段距离处（绘制 50 dp，长按 60 dp），并显示准星与吸附提示。习惯了之后，精度不输鼠标。
\end{tip}

\section{中文界面与语言切换}

中文版会\textbf{自动跟随系统语言}：系统是中文就显示简体中文，否则显示英文。也可以手动固定：

\begin{enumerate}
  \item 打开 \menupath{编辑 → 偏好设置}（或点工具条上的 \ui{偏好设置}）；
  \item 选 \ui{语言（当前：…）}；
  \item 在列表里选 \ui{系统语言}、\ui{英语} 或 \ui{简体中文}。界面会立即重绘，文档与引擎状态保留。
\end{enumerate}

\screenshot[0.8]{ui-language-zh.jpg}{语言选择对话框。第一项“系统语言（简体中文）”表示跟随 Android；切换后界面即时刷新。}

\begin{tip}[术语一致性]
中文界面里的命令名、任务面板字段与属性名，直接采用 FreeCAD 官方中文翻译（例如 \ui{凸台}、\ui{凹槽}、
\ui{完全约束}）。如果你同时使用桌面 FreeCAD，两边会看到同一套术语。
\end{tip}

\section{配置自己的工作区}

在 \menupath{编辑 → 偏好设置} 里还有三项常用设置：

\begin{itemize}
  \item \textbf{单位系统}：切换 FreeCAD 的单位方案（毫米/米、英寸等）。影响所有数值的显示与输入；
  \item \textbf{主题}：跟随系统、浅色或深色。深色适合长时间建模与 OLED 屏幕；
  \item \textbf{每日检查更新}：开启后启动时静默检查 GitHub Releases，有新版会弹出更新对话框（显示版本说明、
        下载大小，并提示“安装时 Android 会关闭 MobileCAD，请先保存”）。
\end{itemize}

\screenshot[0.8]{ui-preferences-zh.jpg}{偏好设置：单位系统、主题、语言、更新检查。}

\chapter{界面详解}

MobileCAD 的界面只有四个常驻区域：\textbf{顶栏}、\textbf{特征树}、\textbf{工具条与视图区}、\textbf{元素选项卡}。
所有对话框、面板、提示都以浮层形式出现在视图区之上。理解这四块，就理解了整个应用。

\screenshot{ui-main-zh.jpg}{主界面的四个常驻区域：顶栏（1）、特征树（2）、工具条与视图区（3）、元素选项卡（4）。}

\section{顶栏}

从左到右：

\begin{itemize}
  \item \ui{☰}：主菜单（文件、编辑、视图、工具、宏、帮助）。
  \item \textbf{文档名}：当前文档（\file{.FCStd}）的名字；未保存时显示“未命名”。
  \item \textbf{副标题}：\ui{工作台 · 基于 FreeCAD 提交号}，例如“零件 · 基于 FreeCAD main e95a31c”。
        它告诉你现在处于哪个工作台、引擎对应的 FreeCAD 版本。
  \item \ui{↶ / ↷}：撤销与重做。长按（或看提示文字）会显示将要撤销的操作名，例如“撤销 凸台”。
\end{itemize}

\section{元素选项卡与工作台}

底部一排选项卡对应 FreeCAD 的工作台，但按“设计流程”重新组织过：

\begin{center}
\small
\begin{tabularx}{\linewidth}{@{}lX@{}}
\toprule
\textbf{选项卡} & \textbf{内容} \\
\midrule
\ui{零件} & 零件设计（PartDesign）+ 零件（Part）+ 草图（Sketcher）+ 曲面（Surface）+ 绘图（Draft）的组合工作区，日常建模都在这里 \\
\ui{装配} & 装配（Assembly）：关节、求解、BOM \\
\ui{工程图} & 工程图（TechDraw）：出二维图并导出 PDF/SVG/DXF \\
\ui{变量} & 电子表格（Spreadsheet）：参数表、别名、表达式 \\
\ui{更多 ▾} & BIM、FEM、CAM、Mesh、Mesh Part、Points、Reverse Engineering、Robot、OpenSCAD、Inspection、Material、Draft、Part、Part Design、Sketcher、Surface 等全部工作台 \\
\bottomrule
\end{tabularx}
\end{center}

切换工作台有两种方式：点底部选项卡；或用菜单 \menupath{视图 → 工作台} 从完整列表中选。切到“更多”里的某个
工作台后，标签会显示其名字（如“FEM ▾”）。

\section{特征树}

左侧是 \textbf{模型树}（FreeCAD 的 Tree View）。它不只是目录，而是模型的“历史与结构”：

\begin{itemize}
  \item \textbf{层级}：实体（Body）包含草图、特征；装配包含构件与关节；工程图包含页面与视图。
  \item \textbf{图标}：每种对象有对应图标（体、凸台、草图、基准面、点、孔、阵列……），与 FreeCAD 桌面一致。
  \item \textbf{状态}：红色文字表示特征出错；灰色/半透明表示隐藏；带删除线表示被压缩（Suppressed）；
        名称后带小图标表示需要重新计算（Touched）；\ui{Tip} 标记表示实体当前计算到哪一步。
      特征行右侧的小红叹号表示该特征失败，打开它的对话框会看到具体错误。
  \item \textbf{单击}：选中对象（在视图中高亮）；\textbf{双击}：打开它的对话框（实体上双击则切换“活动实体”）；
        \textbf{长按}：右键菜单。
  \item \ui{◀}：收起特征树（收起后变成 40 dp 的竖条，点一下再展开）；平板上有 14 dp 拖动条，可把树宽在
        160--480 dp 之间调整（默认 240 dp）；手机上特征树是覆盖式抽屉，从顶栏的树按钮打开，点暗色遮罩关闭。
  \item 树顶的计数（例如“15”）是当前文档的对象数。树\textbf{不支持拖拽重排}——模型结构由特征历史决定，
        想改变先后顺序请用“将 Tip 滚动到此处”和特征对话框里的引用关系。
\end{itemize}

\screenshot[0.72]{tablet-partdesign.jpg}{特征树与 PartDesign 模型。树里能看到 Pad、Groove、Mirrored、Revolution、DatumPlane、Pocket、PolarPattern 等特征的完整建模历史，这就是“参数化”的含义：每一步都可回改。}

长按对象弹出的菜单（随对象类型变化）：

\begin{center}
\small
\begin{tabularx}{\linewidth}{@{}lX@{}}
\toprule
\textbf{菜单项} & \textbf{作用} \\
\midrule
\ui{编辑 …} / \ui{编辑附着} & 打开特征对话框；或只编辑它的附着（草图/基准面挂在哪里）\\
\ui{重命名} & 修改显示名称（不改变内部类型）\\
\ui{隐藏/显示} & 切换可见性（也用于子元素）\\
\ui{压缩 / 取消压缩} & 暂时跳过该特征参与计算，用于排查错误或对比设计 \\
\ui{标记为重新计算} & 强制该对象在下次重算时更新 \\
\ui{重新计算对象} & 立即重算该对象 \\
\ui{将 Tip 滚动到此处} & 把实体的计算终点移到该特征——做“历史回放”式的设计探索 \\
\ui{移入实体} & 把外部的圆角/倒角等特征并入实体历史 \\
\ui{切换活动实体} & 多实体建模时指定新特征进入哪个实体 \\
\ui{删除} & 删除对象（可用撤销恢复）\\
\bottomrule
\end{tabularx}
\end{center}

\section{工具条与工具面板}

工具条上是一排\textbf{工具组}（toolset），例如 \ui{草图}、\ui{特征}、\ui{零件}、\ui{曲面}、\ui{Draft}。
每个工具组对应 FreeCAD 的一个或多个工具栏：

\begin{itemize}
  \item 点工具组 → 在视图区左上角展开\textbf{工具面板}：按 FreeCAD 工具栏分节排列按钮，按钮名即 FreeCAD 命令名；
  \item 面板顶部有 \ui{搜索所有工具}：输入中英文都能搜（例如“pad”或“凸台”），结果按工作台分组；
  \item 面板底部的 \ui{更多工具（N）} 会折叠/展开“该工作台不在工具栏里的命令”，保证所有命令都能触达；
  \item 工具条右侧会显示\textbf{最近使用}的 6 个工具，用竖线分隔，方便重复操作；手机上工具条可横向滚动；
  \item 再点一次工具组（或点视图空白处）关闭面板。
\end{itemize}

\screenshot{ui-toolpanel-zh.jpg}{“特征”工具面板：按 FreeCAD 工具栏分节（建模特征……），命令名与图标都来自 FreeCAD；顶部可搜索全部工具。}

\section{视图区}

\begin{itemize}
  \item \textbf{导航立方体}（右上角，112×108 dp）：六个面（上/下/前/后/右/左）对应标准视图，点面心动画转过去；
        面与面之间的边条、角块给出 26 个方向；四周的 ▲▼◀▶ 每次转 15°，↶ ↷ 转 90°；右上角的 ☰ 小菜单可切换
        \ui{正交/透视}、\ui{等轴测}、\ui{适合全部}。直接拖动立方体可自由旋转。它永远指示当前朝向。
  \item \textbf{选择}：轻触拾取面、边、顶点；已选元素在视图区底部形成一层“选择条”（标签 + ×），点 × 取消。
  \item \textbf{浮动条}：草图状态条（“完全约束 ›”）、草图数值输入、Draft 提示条、装配配合提示、选择过滤器……都
        浮在视图区边缘，不遮挡模型。
  \item \textbf{右键菜单}：视图上\textbf{双指轻点}（或鼠标右键）弹出上下文菜单，内容随当前选择变化：有面时可以
        “草图/拉伸/平面”，有边时可以“圆角/倒角”，有草图时可以“拉伸/旋转/凹槽/编辑草图”。
\end{itemize}

\section{对话框体系}

MobileCAD 有五种交互控件，认识它们就不会迷路：

\begin{description}
  \item[特征对话框（Feature Dialog）] 对应桌面 FreeCAD 的任务面板：44 dp 标题栏（▾ 折叠、✓ 确定、✗ 取消，
        两个按钮分开摆放防止误触），下方错误条，中间是属性行，底部 \ui{更多选项（N）}/\ui{收起选项}。
        属性行按 FreeCAD 各任务面板的顺序排列（例如凸台：轮廓/类型/长度/长度2/……），主行在前，其余折叠。
        行类型包括：数值（数字键盘）、枚举（≤4 项显示为分段按钮，更多则下拉）、布尔（开关）、链接
        （显示“凸台 · 边3, 边7”，点击进入拾取模式，提示 \ui{点选几何}）、颜色（色块 + 11 色调色板 + 自定义）、
        文本/列表/文件、边表（逐条边设半径）、向量/放置（每个分量一行）。\textbf{每改一项都是一次可撤销事务，
        模型实时预览}；✓ 提交，✗ 回滚到打开时的状态。
  \item[任务面板（Task Panel）] 由 FreeCAD 的 Python 命令自己构建的界面（BIM、CAM、FEM、OpenSCAD、TechDraw 等），
        被渲染成同样的触屏控件：分组框、数值框、下拉框、列表、表格、复选/单选、按钮。底部是按 FreeCAD 定义的
        \ui{确定/取消/应用/关闭}。
  \item[选择器（Chooser）] 一个标题 + 若干选项（单选或多选）。用于“选择平面”“选择关节类型”“选择要显示的视图”等。
  \item[数字键盘（Numpad）] 需要数值时从屏幕边缘弹出（平板在右侧 320 dp，手机在底部全宽）：
        键位 7 8 9 / 4 5 6 / 1 2 3 / ± 0 .，右列是退格、+、−、×、÷，底行是括号、C、✕、✓；
        单位可在顶部切换（长度 mm/cm/m/µm/in/ft，角度 °/rad）。输出就是 FreeCAD 表达式，例如 \file{25 mm}、
        \file{2 * (10 mm + 5 mm)}。属性值以 \file{=} 开头时按\textbf{表达式}处理（蓝色显示）。
  \item[文本输入] 单行或多行文本框，用于文件名、标签、表达式、宏代码等。
\end{description}

\screenshot{ui-feature-dialog-zh.jpg}{特征对话框：标题“圆锥体”，属性 角度/高度/半径 1/半径 2，底部“更多选项（6）”。左上角的箭头可折叠，右上角 ✓ 确定、✗ 取消。}

\section{偏好设置}

\menupath{编辑 → 偏好设置}（工具条上也有 \ui{偏好设置}）提供四类设置：单位系统、主题、语言、更新检查。
详见 2.6 节。

\section{平板、手机与头显}

\begin{itemize}
  \item \textbf{平板（宽 ≥ 600 dp）}：特征树常驻左侧，工具面板浮在视图左上角，对话框浮在左侧；
  \item \textbf{手机（宽 < 600 dp）}：特征树变成抽屉（点顶栏树按钮或左滑），面板从底部弹出，工具条横向滚动；
  \item \textbf{MobileCAD XR}（Quest 等头显）：整个界面渲染在一块 880×620 dp 的面板里，模型在房间中，
        数字键盘与软键盘显示为第二块浮窗；对话框改为内嵌绘制（代码中的 LocalInlineDialogs）。
\end{itemize}

\section{返回键、撤销与“卡住了”}

\begin{itemize}
  \item \textbf{返回键就是 Esc}：它按“从内到外”的顺序关闭东西——数字键盘 → 文本对话框 → 选择框 → 菜单 →
        工具面板 → Draft 工具 → 装配配合 → 任务面板 → 特征对话框 → 草图；全都关着时，应用退到后台
        \textbf{而不是退出}，避免误触丢失文档。
  \item \textbf{撤销/重做}在顶栏右侧。每个特征修改、每次属性编辑都是独立事务；撤销按钮的提示会写出具体操作名。
  \item 命令不可用时（菜单置灰、工具面板 38\% 透明），运行时点按会提示“<命令> 尚未在移动端提供”——
        这类命令通常属于桌面窗口管理，移动端没有对应物。
\end{itemize}

\section{键盘与输入法}

\begin{itemize}
  \item 需要输入文字时（表达式、文本属性、宏），软键盘自动弹出；Android 的输入法候选词可用。
  \item 连接\textbf{硬件键盘}时，输入框可直接键入；物理按键不会被输入法截获（应用为此专门处理了 Android 的 IME 焦点）。
  \item 单行输入框按回车 = 确认并收起键盘；多行输入框回车 = 换行。
\end{itemize}

\chapter{第一个零件：二十分钟做出一个支架}

本章是一个完整的、可复现的练习。跟着做，你会用到\textbf{草图 → 凸台 → 圆角 → 孔 → 保存}这条最核心的
建模链路。做完之后，你已经具备用 MobileCAD 做简单零件的能力。

\begin{practice}[练习目标]
一个 60 × 40 mm、厚 12 mm 的长方体支架，四角 R6 圆角，中心一个 ⌀10 的通孔，两侧各一个 ⌀5 的安装孔。
\end{practice}

\section{第 1 步：新建文档}

启动后默认就是“未命名”空文档。若已有打开的文档，用 \menupath{文件 → 新建}（\cmd{Std\_New}）建立新文档。
先保存一次，给文档起个名字：

\begin{enumerate}
  \item \menupath{文件 → 保存}（\cmd{Std\_Save}）；
  \item 在系统文件选择器里选一个目录（例如“文档”），文件名输入 \file{bracket}；
  \item 确定。顶栏的文档名会变成 \file{bracket}。
\end{enumerate}

\begin{tip}[为什么先保存？]
MobileCAD 会定时自动保存，并在崩溃后从恢复副本还原；但“先保存”能确保这次工作有明确的落点，也方便之后
另存副本做设计对比。
\end{tip}

\section{第 2 步：新建草图并选基准面}

\begin{enumerate}
  \item 确认底部在 \ui{零件} 选项卡，点顶部工具条的 \ui{草图}；
  \item 在工具面板里点 \ui{新建草图}（命令 \cmd{Sketcher\_NewSketch}）；
  \item 弹出 \ui{草图平面} 选择：\ui{上 (XY)}、\ui{前 (XZ)}、\ui{右 (YZ)}。选 \ui{上 (XY)}。
\end{enumerate}

\screenshot{ui-sketch-plane-zh.jpg}{新建草图时选择基准面。之后可以在草图属性里用“附着”改挂到实体表面或其他基准面上。}

\section{第 3 步：画矩形}

进入草图模式后，工具条会变成草图专用：\ui{绘制}、\ui{约束}、\ui{修改}、\ui{视图} 四个工具组，
以及 \ui{构造：关}、\ui{约束}（约束列表）、常用的 \ui{折线}、\ui{矩形}、\ui{圆}、\ui{三点圆弧}、
\ui{尺寸}、\ui{修剪}、\ui{选择}、\ui{删除}，最右侧是 ✓（离开草图）与 ✗（取消草图）。

\begin{enumerate}
  \item 点 \ui{矩形}；
  \item 在原点附近轻触一次（第一角点），移动手指，再轻触一次（对角点）。
        不用追求精确尺寸，形状对就行——下一步用尺寸“锁死”它。
\end{enumerate}

\screenshot{ui-sketch-tools-zh.jpg}{草图工具面板：“草图”工具组按几何、约束、编辑、视图分节列出全部草图命令；进入草图后，顶部工具条会切成草图专用的绘制/约束/修改/视图四组。}

\section{第 4 步：标注尺寸，让草图“完全约束”}

\begin{enumerate}
  \item 点 \ui{尺寸}（\cmd{Sketcher\_Dimension}），依次点矩形的两条水平边或两个角点，拖到合适位置轻触放置；
  \item 放置后数字键盘弹出，输入 \file{60}，✓ 确认；
  \item 同样标注另一方向，输入 \file{40}；
  \item 观察左下角状态：当显示 \ui{完全约束}（Fully constrained）时，草图的每个自由度都被约束住了；
        若显示“欠约束 N 个自由度”，说明还差约束，模型仍可能被意外拖动。
\end{enumerate}

\begin{tip}[完全约束为什么重要？]
完全约束意味着几何形状被数学上唯一确定。后续特征（凸台、凹槽）依赖这张草图：欠约束的草图在参数变化时
可能“跳形”，导致特征失败。养成习惯：\textbf{每次离开草图前，让状态条变成“完全约束”}。
\end{tip}

\section{第 5 步：离开草图}

点工具条最右侧的 ✓（提示“离开草图”）。视图回到 3D，特征树里多出一个 \ui{草图}（Sketch），它被收在新建的
\ui{实体}（Body）里——这是 PartDesign 的规矩：所有特征都属于某个实体。

\section{第 6 步：凸台（Pad）}

\begin{enumerate}
  \item 在树里点选刚画的草图；
  \item 点工具条的 \ui{特征} → \ui{凸台}（\cmd{PartDesign\_Pad}）；
  \item 特征对话框出现，\ui{轮廓} 已经自动填好；在 \ui{长度} 行点一下，用数字键盘输入 \file{12}；
  \item 想直接看到拉伸方向？在 \ui{反向} 开关上切换；
  \item 点标题栏的 ✓ 确认。长方体出现了。
\end{enumerate}

\begin{warn}[特征失败怎么办？]
如果对话框下方出现红色错误（例如“无法在基准特征上找到该几何”），说明轮廓有问题（未闭合、自相交、多余线）。
取消后回到草图，用 \ui{约束} 列表检查冲突，用 \ui{修剪} 清掉多余的线段，再试一次。
\end{warn}

\section{第 7 步：四角圆角}

\begin{enumerate}
  \item 点工具条的 \ui{特征} → \ui{圆角}（\cmd{PartDesign\_Fillet}）；
  \item 对话框里的 \ui{基体} 已经指向当前实体；点 \ui{边}（或直接在模型上轻触四条竖直边）；
  \item 在 \ui{半径} 里输入 \file{6}；
  \item ✓ 确认。四条竖边变成 R6 圆角。
\end{enumerate}

\begin{tip}[圆角和倒角的“边表”]
圆角/倒角对话框可以逐条边设置不同半径：每选一条边就多一行，点行里的数值可单独修改。想一次全部改，
用 \ui{全部半径} 按钮。
\end{tip}

\section{第 8 步：中心通孔}

有两种做法，推荐第一种：

\subsection{方法 A：孔特征（Hole）}

\begin{enumerate}
  \item 先选\textbf{顶面}（轻触上表面）；
  \item 点 \ui{特征} → \ui{孔}（\cmd{PartDesign\_Hole}）；
  \item 在对话框里设置 \ui{直径} 为 \file{10}，\ui{深度类型} 为 \ui{贯穿全部}（Through all）；
  \item ✓ 确认。
\end{enumerate}

\begin{tip}[标准螺纹与沉头]
孔特征的 \ui{螺纹} 开关会按 ISO 标准生成螺纹孔（选螺纹类型、规格、精度）；\ui{孔类型} 可选沉头、埋头、
沉孔等，并自动显示对应的直径、深度、角度参数。
\end{tip}

\subsection{方法 B：草图 + 凹槽（Pocket）}

在顶面上新建草图 → 画一个 ⌀10 的圆 → 离开草图 → 选草图点 \ui{凹槽}（\cmd{PartDesign\_Pocket}）→
深度选“贯穿全部”。两种方法结果相同，方法 B 更灵活（可以画异形孔）。

\section{第 9 步：两个安装孔（练习）}

用同样的方法，在两侧各开一个 ⌀5 的通孔。提示：
\begin{itemize}
  \item 用 \ui{圆} 工具画圆，加 \ui{尺寸} 约束直径（输入 \file{5}），再用 \ui{距离} 约束定位；
  \item 也可以先画一个孔，再用 \ui{极阵列} 或 \ui{线性阵列} 复制；
  \item 完成后树里应该依次是：草图 → 凸台 → 圆角 → 孔 → 草图 → 凹槽 → 圆角……
\end{itemize}

\section{第 10 步：回头改参数}

这才是参数化 CAD 的精髓：

\begin{enumerate}
  \item 双击特征树里的 \ui{凸台} → 把长度从 12 改成 20 → ✓，整个模型连同后面的圆角、孔一起更新；
  \item 双击第一个 \ui{草图} → 把 60 改成 80 → 离开草图，凸台自动变长；
  \item 若某个后续特征因为尺寸变化而失败（例如孔跑到了实体外），树里会出现红色错误，双击那个特征看原因，
        改约束或改位置即可。
\end{enumerate}

\section{第 11 步：保存与导出}

\begin{itemize}
  \item \menupath{文件 → 保存}：保存为 \file{.FCStd}（FreeCAD 原生格式，保留全部历史）；
  \item \menupath{文件 → 另存副本}：存一份快照；
  \item \menupath{文件 → 导出}：可选 STEP、IGES、STL、OBJ、3MF 等，供 3D 打印或其它 CAD 使用。
\end{itemize}

\begin{tip}[给 3D 打印的建议]
导出 STL 前先用 \ui{零件 → 检查几何} 检查实体是否有效；STL 的精度取决于“偏差”参数，平板上的手办类模型
建议 0.05 mm 左右，机械件 0.01 mm。
\end{tip}

\section{小结}

你已经走完了 CAD 的主干流程：

\begin{center}
\small
\begin{tabularx}{\linewidth}{@{}lX@{}}
\toprule
\textbf{概念} & \textbf{在 MobileCAD 里对应什么} \\
\midrule
草图 & 2D 轮廓 + 约束 + 尺寸；完全约束后才稳定 \\
凸台 / 凹槽 & 把闭合轮廓沿法向加料 / 减料 \\
圆角 / 倒角 & 对边做修饰，位于特征历史中 \\
孔 & 标准化开孔，支持螺纹、沉头 \\
特征树 & 建模历史；双击任何一步都能回改 \\
Tip & 实体“算到哪一步”的标记，可滚动回放 \\
\bottomrule
\end{tabularx}
\end{center}

下一章我们把草图讲透：所有绘制工具、所有约束类型、冲突排查与 B 样条。

\chapter{草图：一切参数化建模的起点}

草图（Sketcher）是 PartDesign 的地基：\textbf{2D 轮廓 + 约束 + 尺寸}。本章把草图工作台拆到底——
每个绘制工具需要点几次、每个约束怎么用、冲突怎么排查、B 样条怎么编辑。

\section{进入与退出草图}

\subsection{新建草图}

\begin{enumerate}
  \item 点工具条的 \ui{草图}（Sketcher 工具组）→ \ui{新建草图}（\cmd{Sketcher\_NewSketch}）；
  \item 若当前已选中一个平面或基准面，草图直接附着到它上面；否则弹出 \ui{草图平面} 选择：
        \ui{上 (XY)}、\ui{前 (XZ)}、\ui{右 (YZ)}；
  \item 草图会落在\textbf{活动实体}（Body）中，附着模式为“贴平面”（FlatFace）。曲面不能作为草图平面。
\end{enumerate}

\subsection{编辑与退出}

\begin{itemize}
  \item 编辑已有草图：树里双击草图，或选中后点 \ui{编辑草图}（\cmd{Sketcher\_EditSketch}）；
  \item \ui{✓ 离开草图}：保存修改并重算（快捷键位置在工具条最右）；
  \item \ui{✗ 取消草图}：\textbf{回滚}到打开草图时的状态。两个按钮之间隔 20 dp，防止误触。
\end{itemize}

\begin{warn}[撤销在草图中是“封底”的]
进入草图后，撤销只能退到“打开草图”这一步之前——草图内部的所有修改属于同一个事务。想整段放弃，
请用 ✗ 取消草图；取消掉的草图仍然可以用“重做”恢复。
\end{warn}

\section{草图工具条}

进入草图后，顶部工具条变成四组：

\begin{center}
\small
\begin{tabularx}{\linewidth}{@{}lX@{}}
\toprule
\textbf{工具组} & \textbf{内容} \\
\midrule
\ui{绘制} & 折线、线、矩形、圆、圆弧等全部绘制工具（对应 FreeCAD 的 Geometries 工具条）\\
\ui{约束} & 全部约束命令（对应 Constraints 工具条）\\
\ui{修改} & 修剪、延伸、分割、圆角、倒角、移动、旋转、缩放、阵列、外部几何、B 样条工具 \\
\ui{视图} & 对齐视图、剖切、网格、显示顺序、圆弧叠加、B 样条信息层 \\
\bottomrule
\end{tabularx}
\end{center}

工具组右侧还有固定按钮：\ui{构造：开/关}、\ui{约束}（打开约束列表）、\ui{折线}、\ui{矩形}、\ui{圆}、
\ui{三点圆弧}、\ui{尺寸}、\ui{修剪}、\ui{选择}、\ui{删除}、\ui{完成}，以及最右的 ✓ / ✗。
工具面板里同样按 FreeCAD 工具条分节，并有 \ui{更多工具（N）}。

\screenshot{ui-sketch-tools-zh.jpg}{草图模式：顶部工具条四组 + 固定按钮；底部状态条显示“完全约束”；右下角是当前工具（折线）的退出胶囊。}

\section{绘制工具全表}

每个绘制工具需要“点”若干次；多点工具（折线、B 样条）用\textbf{双击}或点 \ui{完成} 结束。

\begin{longtable}{@{}p{0.20\linewidth}p{0.12\linewidth}p{0.60\linewidth}@{}}
\toprule \textbf{工具} & \textbf{点数} & \textbf{操作} \\ \midrule \endhead
线 & 2 & 起点 → 终点 \\
折线 & 多点 & 逐段点；点回起点闭合；双击或“完成”结束 \\
矩形 & 2 & 两个对角点 \\
中心矩形 & 2 & 中心 → 角点（自动加构造线与对称）\\
圆 & 2 & 圆心 → 半径点 \\
三点圆 & 3 & 圆上三点 \\
圆弧 & 3 & 圆心 → 起点 → 终点（逆时针）\\
三点圆弧 & 3 & 起点 → 终点 → 弧上一点 \\
点 & 1 & 放置一个点 \\
长圆槽 & 3 & 两圆心 → 半径点 \\
正多边形 & 2 & 中心 → 角点（边数由命令定：3--8 边）\\
椭圆 & 3 & 中心 → 长轴端 → 定点 \\
三点椭圆 & 3 & 椭圆上三点 \\
椭圆弧 & 4 & 中心、长轴、起点、终点 \\
双曲线弧 & 4 & 焦点、顶点、起点、终点 \\
抛物线弧 & 4 & 焦点、顶点、起点、终点 \\
圆角矩形 & 2 & 两个对角点（带圆角）\\
圆弧槽 & 4 & 弧上两点 + 两个半径点 \\
文本 & 2 & 先画一条把手线，再输入文字 \\
B 样条 & 多点 & 控制点式 / 周期式 / 按节点插值式；双击或“完成”结束 \\
\bottomrule
\end{longtable}

\subsection{动作型工具（不画新几何）}

修剪、尺寸、圆角、倒角、延伸、分割、连接曲线、对称、移动、复制、旋转、缩放、矩形阵列、外部投影/相交、
插节点、增减重数、粘贴——这些工具都是“点选已有几何后执行”，详见 5.10 与 5.12。

\section{吸附与自动约束}

草图里有一个总开关 \cmd{Sketcher\_Snap}（工具面板的 \ui{吸附/自动约束}）。开启时，落点会自动吸附，
并\textbf{自动加上合适的约束}：

\begin{itemize}
  \item 落在端点、圆心、原点 → \textbf{重合}（Coincident）；
  \item 落在模型顶点/圆心 → 重合，并把该元素隐式投影进草图；
  \item 落在两曲线交点 → 两条“点在对象上”；
  \item 落在线段中点 → 自动加“关于两端点对称”；落在圆弧中点 → “点在对象上”；
  \item 落在曲线或坐标轴上 → “点在对象上”；
  \item 与上一点接近水平/竖直（2° 容差）→ 自动加水平/竖直约束，并显示绿色对齐虚线。
\end{itemize}

光标旁会直接写出将添加的约束（例如“coincident + horizontal”），加完还会自动删除冗余约束。关掉吸附＝同时
关掉推断；网格由 \cmd{Sketcher\_Grid} 单独控制。

\begin{tip}[自动约束是双刃剑]
它让画图飞快，但也可能在你没注意时加上“不想 要”的约束，导致后面改尺寸跳形。养成习惯：画完后打开
\ui{约束} 列表过一遍，把多余的重合/水平竖直删掉。
\end{tip}

\section{约束类型全表}

\begin{center}
\small
\begin{tabularx}{\linewidth}{@{}lX@{}}
\toprule
\textbf{约束} & \textbf{用法与效果} \\
\midrule
重合 Coincident & 两点（或点与曲线端点）重合 \\
点在对象上 PointOnObject & 点落在曲线/轴上 \\
水平/竖直 Horizontal/Vertical & 线水平或竖直；\cmd{HorVer} 会按当前方向自动判断 \\
平行 Parallel / 垂直 Perpendicular & 两条线的关系 \\
相切 Tangent & 两曲线相切，或曲线与端点相接处的相切 \\
相等 Equal & 等长、等半径 \\
对称 Symmetric & 两点关于点/线对称 \\
锁定 Block & 固定已有几何，使其不被后续操作拖动 \\
固定 Lock & 把点同时约束到原点的 X、Y 距离上 \\
距离 Distance / 距离X / 距离Y & 长度或坐标分量 \\
半径 Radius / 直径 Diameter & 圆/圆弧尺寸；\cmd{Radiam} 一键切换 \\
角度 Angle & 两线夹角，或单线相对 X 轴的角度 \\
组 Group & 把若干几何作为一个整体（配合块使用）\\
Snell 定律 SnellsLaw & 光学折射约束 \\
\bottomrule
\end{tabularx}
\end{center}

\section{施加约束的两种方式}

\begin{enumerate}
  \item \textbf{先选元素，再点约束}：点选一个/多个几何（点、线、圆），再点约束命令，立即生效。
  \item \textbf{先点约束，再点元素（武装模式）}：什么都不选就点约束命令，状态条会提示依次点哪些元素；
        点够数量自动施加，并可连续施加；若点的元素类型不匹配，会提示 \ui{Not a valid selection for X}
        并清空选择。
\end{enumerate}

\section{尺寸：一把统一的“尺寸”命令}

\cmd{Sketcher\_Dimension} 会按当前选择自动判断要标什么：长度、半径、直径或角度，并弹出数字键盘、预填
当前值。X/Y 距离会自动交换两点使结果为正值；也支持“点到线的垂距”。

\begin{itemize}
  \item \textbf{修改尺寸}：长按尺寸标签（触感反馈）→ 数字键盘；或打开约束列表点行末的 ✎；也可直接拖动标签位置；
  \item \textbf{驱动 / 参考}：列表行菜单里的 \ui{转为参考（从动）} 让尺寸只测量不参与求解（灰色，标“(reference)”）；
  \item \textbf{命名约束}：给它起名后，可以在表达式里引用，例如 \file{=Sketch.Constraints.Width * 2}。
\end{itemize}

\section{约束列表与冲突排查}

点工具条的 \ui{约束} 打开列表（标题显示约束总数）：

\begin{itemize}
  \item 顶部筛选 \ui{全部} / \ui{冲突 / 冗余}；
  \item 每行显示编号、名称、值、关系描述（如“线3 ↔ 圆5”）；
  \item 颜色：\textbf{红色 = 冲突}、\textbf{橙色 = 冗余}、\textbf{灰色 = 停用}（标“(off)”）；
  \item 行菜单：\ui{停用/启用}、\ui{转为参考（从动）}/\ui{转为驱动}、\ui{重命名}、\ui{删除}。
\end{itemize}

\begin{practice}[冲突处理的标准流程]
\begin{enumerate}
  \item 状态条点开约束列表，筛选 \ui{冲突 / 冗余}；
  \item 对红色冲突：停用其中一条看是否恢复，再决定删除哪条；
  \item 对橙色冗余：直接删除（它只是重复）；
  \item 运行 \ui{验证草图}（Validate Sketch）：它会补重合点、补水平/竖直、自动去除冗余；
  \item 回到 \ui{全部}，确认状态条变回“完全约束”。
\end{enumerate}
\end{practice}

\section{构造几何}

构造几何只用于辅助作图，不参与成形（显示为青色虚线）。

\begin{itemize}
  \item \textbf{没有选中任何曲线}时，点 \ui{构造：关} 切换为“下一个画出的形状是构造线”；
  \item \textbf{选中曲线}后，同一个开关把选中曲线在“构造/普通”之间转换；
  \item 外部几何（投影进来的边）也可以在两态之间切换“参考/定义”。
\end{itemize}

\section{编辑工具}

\begin{center}
\small
\begin{tabularx}{\linewidth}{@{}lX@{}}
\toprule
\textbf{工具} & \textbf{操作要点} \\
\midrule
修剪 Trim & 点要去掉的线段，一次一段 \\
延伸 Extend & 点靠近待延伸端，再输入增量（默认 10 mm）\\
分割 Split & 点分割位置，把一条曲线拆成两条 \\
连接曲线 Join & 依次点两条曲线的端点 \\
草图圆角 / 倒角 & 点一个\textbf{顶点}，输入半径/尺寸（默认 2 mm）\\
移动 / 复制 & 先选中几何，点两次定义位移矢量 \\
旋转 & 选几何 → 点中心 → 输入角度（默认 90°）\\
缩放 & 选几何 → 点中心 → 输入倍率（默认 2）\\
对称 & 选几何 → 点对称线或对称点 \\
矩形阵列 & 选几何 → 点两次定方向与间距 → 输入行列数（默认 2 × 2）\\
偏移 Offset & 选几何后点工具，正数外扩、负数内缩 \\
粘贴 & 先复制/剪切，再用粘贴工具点目标点 \\
\bottomrule
\end{tabularx}
\end{center}

\begin{warn}[这些工具都有前提]
对称、移动、旋转、缩放、阵列、偏移都要求\textbf{先选中几何}，否则只会提示“请先选择几何”。
圆角/倒角只作用于顶点，不是任意两条曲线；修剪/分割一次只处理一条曲线。
\end{warn}

\section{外部几何：引用模型上的边}

\ui{外部投影}（External geometry）把模型上的边/顶点投影进草图，作为绘制参照；\ui{外部相交} 则取与草图
平面相交的面/边。

\begin{itemize}
  \item 选工具后，点模型上的边或顶点；
  \item 跨 Body 引用时，应用会自动改用同一实体内的元素，或创建一个 SubShapeBinder；会造成循环引用时直接报错；
  \item 投影进来的几何默认是“参考”（不参与成形），可切换为“定义”（参与）。
\end{itemize}

\section{B 样条与曲线编辑}

四种创建方式：\textbf{控制点式}、\textbf{周期式}（至少 3 点）、\textbf{按节点插值式}、\textbf{周期插值式}。
度数取 $\min(3, n-1)$。

\begin{itemize}
  \item \ui{插入节点}：点曲线上的位置，在最近点处插入；
  \item \ui{增加/减少节点重数}：必须点在节点上，否则提示 \ui{请选择 B 样条节点}；
  \item \ui{转换为 NURBS}、\ui{增减次数}、\ui{连接曲线} 是曲线级操作；
  \item 视图工具组里的 \ui{B 样条信息} 可分别显示：次数、控制多边形、曲率梳、节点重数、极点权重。
\end{itemize}

\section{草图里的数值输入}

绘制过程中，视图底部会出现 \ui{输入：} 行，可为主题形状键入精确值：

\begin{center}
\small
\begin{tabularx}{\linewidth}{@{}lX@{}}
\toprule
\textbf{形状} & \textbf{可输入} \\
\midrule
线 / 折线 & 长度、角度 \\
矩形 / 中心矩形 & 宽、高 \\
圆 & 直径 \\
圆弧 / 正多边形 & 半径 \\
\bottomrule
\end{tabularx}
\end{center}

方向取手指最后的位置。槽、椭圆、B 样条等没有输入框，先画后标尺寸。输入的值会变成 FreeCAD 表达式：
纯数字+单位按数值写入，其余（带运算）按表达式写入。

\section{草图视图辅助}

\ui{视图} 工具组里：\ui{对齐视图到草图}（正对草图）、\ui{剖切视图}、\ui{网格}、\ui{显示顺序}、
\ui{圆弧叠加}、\ui{B 样条信息}。网格间距与显示由 FreeCAD 的草图网格设置控制。

\begin{practice}[练习：完全约束的六边形]
\begin{enumerate}
  \item 新建草图，选 XY 平面；
  \item 点 \ui{多边形}（6 边），中心放在原点附近，拉出外接圆；
  \item 加 \ui{尺寸}：外接圆半径（例如 20 mm）；再给中心点加 \ui{固定}（到原点的 X/Y 距离为 0）；
  \item 观察状态条：显示“完全约束”；六条边应全部变黑（已定义）；
  \item 试着拖动任意一个顶点——完全约束时它不会动；
  \item 长按半径尺寸改成 30，六边形整体放大。
\end{enumerate}
\end{practice}

\section{常见问题}

\begin{description}
  \item[草图状态一直是“欠约束 N 个自由度”] 缺 N 个约束。常见：没有固定位置（缺 X/Y 距离）、缺水平/竖直、缺尺寸。
  \item[改尺寸后形状翻转] 草图欠约束时可能整体镜像。加到完全约束，并给关键点加“锁定”。
  \item[多点工具画完却没生成] 折线和 B 样条需要双击或点“完成”提交。
  \item[跨 Body 引用报错] 草图只能在一个 Body 内；需要引用别的体，先建 SubShapeBinder 或改用外部几何。
  \item[想撤销整个草图] 用 ✗ 取消草图，而不是连续撤销。
\end{description}

下一章：把草图变成实体——零件设计工作台。

\chapter{零件设计：把草图变成实体}

零件设计（PartDesign）是 MobileCAD 最常用的工作台。它的核心思想是\textbf{历史}：一个实体（Body）由一串
特征组成，每个特征都是“在上一步结果上做一件事”。改任何一步，后面的步骤自动重算。

\section{三个关键概念：实体、特征、Tip}

\begin{description}
  \item[实体（Body）] 特征的容器，也是 PartDesign 的工作单位。所有草图与特征都属于某个实体。
        新建特征时，如果还没有实体，应用会自动创建一个并把新特征放进去。
  \item[特征（Feature）] 历史中的一步：凸台、凹槽、圆角、阵列……
  \item[Tip] 实体“计算到哪一步”的标记。树里有一条蓝色回滚条，长按某个特征选 \ui{将 Tip 滚动到此处}
        就能回到那一步的历史状态（回放建模过程）。
\end{description}

\textbf{活动实体}（Active Body）：有多个实体时，新特征进入“活动”的那个（树里标签加粗）。双击实体或
右键 \ui{切换活动实体} 来切换。

\section{特征是怎么“被准备好”的}

PartDesign 的命令大多“看选区行事”，这能省掉大量对话框操作：

\begin{itemize}
  \item 选中\textbf{草图}后点凸台/凹槽 → 轮廓自动填入；
  \item 选中\textbf{平面}后点凸台/凹槽 → 先在该平面新建草图；
  \item 什么都没选时，凸台/凹槽会取文档里\textbf{最后一个草图}；没有草图则提示“请先选择草图或平面”；
  \item 圆角/倒角/拔模/厚度：若选中整个特征，则自动对所有边/面生效（\ui{UseAllEdges}）；没选时在 Tip 上等点边；
  \item 默认值：凸台 10 mm、凹槽 5 mm、旋转 360°（轴取草图 V 轴）、镜像默认 V 轴、
        线性阵列 X 向 100 mm × 2 个、极阵列 Z 轴 360° × 2 个。
\end{itemize}

\section{加料特征}

\begin{longtable}{@{}p{0.17\linewidth}p{0.10\linewidth}p{0.65\linewidth}@{}}
\toprule \textbf{特征} & \textbf{命令} & \textbf{关键参数} \\ \midrule \endhead
凸台 & Pad & 轮廓、类型（长度/双向长度/到头/到面/到形状/尺寸）、长度、长度2、偏移、中面、反向、拔模角、精炼 \\
旋转 & Revolution & 轮廓、参考轴、类型（角度/两角度/到头/到面）、角度、角度2、中面、反向、精炼 \\
放样 & Additive Loft & 轮廓、截面（多个）、直纹、闭合、精炼 \\
管道 / 扫掠 & Additive Pipe & 轮廓、路径、模式、过渡、变换、精炼 \\
螺旋 & Additive Helix & 轮廓、参考轴、模式、螺距、高度、圈数、角度、增长、左旋、反向、精炼 \\
加料基本体 & Additive Box/Cylinder/Sphere/Cone/Torus/Ellipsoid/Prism/Wedge & 长宽高/半径/角度等；选中面时贴面创建 \\
\bottomrule
\end{longtable}

\section{减料特征}

\begin{itemize}
  \item \ui{凹槽}（Pocket）：与凸台对称，把轮廓向内挖；深度类型同样支持贯穿全部（Through all）；
  \item \ui{孔}（Hole）：标准孔特征。参数包括：轮廓、\ui{螺纹}、螺纹类型/规格/等级/配合、直径、方向、
        孔切割类型（沉孔/埋头/沉头）及其直径/深度/角度、深度类型/深度、钻尖角度、锥度、反向；
  \item \ui{槽}（Groove）：绕轴挖环形槽；
  \item \ui{减料放样/管道/螺旋} 与加料版本参数一致，只是方向相反；
  \item \ui{减料基本体}：8 种基本体的减料版本。
\end{itemize}

\begin{tip}[孔特征会自动找位置]
在没有圆孔的面直接调用“孔”，应用会在该面中心自动建一个 R3 的草图圆作为轮廓——之后再改孔直径即可。
这能省掉“先画圆草图”的步骤。
\end{tip}

\section{修饰特征}

\begin{description}
  \item[圆角 Fillet] 参数：基体、边（边表，可逐边设半径）、全部使用边、精炼。拖动箭头能直接在模型上改半径。
  \item[倒角 Chamfer] 参数：基体、倒角类型（等距/两距离/距离+角度）、尺寸、尺寸2、角度、翻转方向、全部使用边。
  \item[拔模 Draft] 参数：基体、中性面、拉动方向、角度、反向。
  \item[厚度 Thickness] 参数：基体、厚度值、模式（皮肤/管道/矩形）、连接类型（弧/切线/相交）、反向、相交。
  \item[去除特征] 删除已有的面/圆角，恢复其下的几何。
\end{description}

\screenshot{ui-feature-dialog-zh.jpg}{特征对话框（圆锥体）：主属性在前（角度、高度、半径 1、半径 2），其余收进“更多选项（6）”。右上角 ✓/✗ 决定提交或回滚。}

\section{阵列与镜像}

\begin{center}
\small
\begin{tabularx}{\linewidth}{@{}lX@{}}
\toprule
\textbf{特征} & \textbf{参数} \\
\midrule
线性阵列 LinearPattern & 原始特征、方向（边/轴）、模式（长度/间隔）、长度、偏移、数量、反向、精炼 \\
极阵列 PolarPattern & 原始特征、轴、模式（角度/间隔）、角度、偏移、数量、反向、精炼 \\
镜像 Mirrored & 原始特征、镜像平面、精炼 \\
多变换 MultiTransform & 组合多个阵列/镜像（走通用建对象流程）\\
\bottomrule
\end{tabularx}
\end{center}

\begin{tip}[阵列的“长度 vs 间隔”]
\textbf{长度}（Extent）是首尾之间的总距离；\textbf{间隔}（Spacing）是相邻两个实例的间距。
六个圆周孔：选 \ui{极阵列}，轴取 Z 轴，角度 360°，数量 6 即可。
\end{tip}

\section{参考几何与附着}

\ui{参考几何}（Datum）让你在没有现成面可用时有“虚拟基准”：

\begin{itemize}
  \item \ui{参考平面}、\ui{参考线}、\ui{参考点}、\ui{坐标系}；
  \item 创建时的参数包括：附着支持（\ui{AttachmentSupport}）、贴图模式（\ui{MapMode}）、\ui{反向}、
        偏移（\ui{AttachmentOffset}）、路径参数（\ui{MapPathParameter}）；
  \item \ui{编辑附着}（树里长按）只改附着，不打开整个特征。
\end{itemize}

\begin{warn}[附着模式决定“跟着谁动”]
草图挂在实体表面上时，如果那个面因为前面的特征变化而移动，草图和它派生出的特征会一起更新——这正是
参数化建模想要的行为；但如果面消失（例如被倒角切掉），依赖它的特征就会失败。
\end{warn}

\section{布尔操作}

\begin{itemize}
  \item \ui{零件设计 → 布尔}（PartDesign Boolean）：把选中的对象作为工具，与活动实体做并/差/交；
  \item \ui{零件}（Part）工作台的 \ui{并/差/交} 是“独立实体之间的布尔”，不是 PartDesign 历史的一部分；
  \item 布尔前先确认两个形状都是有效实体（\ui{零件 → 检查几何}）。
\end{itemize}

\section{特征对话框的读法}

特征对话框由三部分组成：

\begin{enumerate}
  \item \textbf{标题栏}：名称 + ▾ 折叠 + ✓ 确定 + ✗ 取消；
  \item \textbf{错误条}（红色）：特征计算失败时显示原因，例如“基体特征上不存在该几何”；
  \item \textbf{属性区}：按 FreeCAD 任务面板顺序排列主属性，其余在 \ui{更多选项（N）} 里。
\end{enumerate}

每种属性行的交互方式见 3.7 节。两个实用细节：

\begin{itemize}
  \item \textbf{实时预览}：改动立即作用于模型，每次改动都是独立事务；
  \item \textbf{拖动箭头}：凸台长度、圆角半径等数量属性在视图里有箭头，可直接拖动并在标签上看到毫米值。
\end{itemize}

\section{失败特征与排查}

特征失败时：

\begin{itemize}
  \item 树里该特征标签变\textbf{红}，并出现红色错误图标；
  \item 视图里用\textbf{红色半透明“幽灵”}显示它最后一次有效的前驱形状；
  \item 打开对话框，顶部红条给出错误字符串。
\end{itemize}

\begin{practice}[排查四步法]
1) 打开失败特征看红条；2) 检查它引用的轮廓/面/边是否还存在；3) 长按它选 \ui{压缩}，看后面的特征是否也依赖它；
4) 回到草图修约束，或长按选 \ui{重新计算对象}。
\end{practice}

\section{完整工作流：从块体到成品}

用第 4 章的支架继续：

\begin{enumerate}
  \item \textbf{加筋}：在支架上表面新建草图，画一条直线 → 退出 → \ui{凸台}（勾选“对称于面”），得到筋板；
  \item \textbf{打沉头孔}：选顶面 → \ui{孔} → 打开 \ui{螺纹} 或选 \ui{沉头} → 填直径/深度；
  \item \textbf{阵列}：选中孔 → \ui{线性阵列}，方向选一条边，数量 4，长度 40；
  \item \textbf{倒角}：选外轮廓边 → \ui{倒角}，尺寸 1 mm；
  \item \textbf{改基体}：双击第一个凸台，把长度从 12 改到 16——筋、孔、阵列、倒角全部自动跟随；
  \item 导出 STEP 给同事，或去“工程图”选项卡出图。
\end{enumerate}

\screenshot[0.72]{tablet-partdesign.jpg}{一个真实的 PartDesign 模型（螺丝刀头）：树里是 Pad → Groove → Mirrored → Revolution → DatumPlane → Pocket → PolarPattern…… 每一步都可以回改，这正是参数化的威力。}

\section{常见坑}

\begin{itemize}
  \item \textbf{草图必须在活动实体内}；跨体引用会被替换成 SubShapeBinder，可能引发循环报错；
  \item \textbf{没有 M 键}：桌面端的折线模式切换、尺寸类型轮换等快捷键在移动端不存在；
  \item \textbf{正多边形边数固定}：命令只有 3--8 边；
  \item \textbf{多点工具要双击/完成}：折线、B 样条不双击不会落库；
  \item \textbf{部分阵列走通用流程}：圆周/路径/点阵列与多变换、椭球/棱柱/楔形的主参数在“更多选项”里；
  \item \textbf{验证草图不显示修复计数}：它确实在修，但界面不回报数字。
\end{itemize}

\chapter{零件工作台：实体、布尔与造型}

\textbf{零件}（Part）工作台与“零件设计”不同：它不维护特征历史，而是直接对\textbf{形状}做操作——建基本体、
做布尔、拉伸/旋转/放样/扫掠、倒圆角。它适合：快速建毛坯、处理导入的 STEP/IGES、做“一次性”几何运算，
以及把多个实体合成为一个。

\section{入口}

底部 \ui{零件} 选项卡的工具条里有五组：\ui{草图}、\ui{特征}、\ui{零件}、\ui{曲面}、\ui{Draft}。
其中 \ui{零件} 展开后分三段：\ui{实体}、\ui{零件工具}、\ui{布尔工具}。也可以从
\menupath{更多 ▾ → 零件} 进入同一工作台。

\section{基本体}

\begin{center}
\small
\begin{tabularx}{\linewidth}{@{}lX@{}}
\toprule
\textbf{基本体} & \textbf{默认参数} \\
\midrule
立方体 Box & 10 × 10 × 10 mm \\
圆柱 Cylinder & 半径 2、高 10、角度 360° \\
球 Sphere & 半径 5、纬度 $-90°\sim90°$、经度 360° \\
圆锥 Cone & 半径1 2、半径2 4、高 10 \\
圆环 Torus & 半径1 10、半径2 2、角度范围 \\
管 Tube & 外半径、内半径、高 \\
\bottomrule
\end{tabularx}
\end{center}

椭球、棱柱、楔形等 16 种图元在 \ui{图元...}（\cmd{Part\_Primitives}）对话框里统一创建：
平面、立方体、圆柱、圆锥、球、椭球、圆环、棱柱（默认 6 边）、楔形、螺旋、螺线、圆、椭圆、顶点、线、
正多边形。

\begin{tip}[参数随时可改]
基本体创建后，双击树里对象即可重新打开参数对话框。尺寸、角度、半径都能改，但\textbf{不会}像 PartDesign
那样带动后续特征——Part 对象之间是独立形状。
\end{tip}

\section{布尔运算}

\begin{itemize}
  \item \ui{布尔}（\cmd{Part\_Boolean}）是对话框：先选两个以上对象，再选运算——\ui{交集 / 剪切 / 结合 / 剖面}；
  \item \ui{并}（Fuse）、\ui{交}（Common）支持多对象（全部选入）；
  \item \ui{差}（Cut）只取“基体 + 工具”：\textbf{第一个选中的是基体}，其余为工具；
  \item “布尔工具”里还有一组变体：\ui{分割形状}（布尔碎片/切片/分离切片/异或）、
        \ui{连接形状}（连接/嵌入/切口）。
\end{itemize}

\begin{warn}[选择顺序决定结果]
Part 的布尔没有“撤销选择顺序”的按钮。切/差之前，养成先点基体、再点工具的习惯；搞反了，撤销重来即可。
\end{warn}

\section{造型：拉伸、旋转、放样、扫掠}

\begin{description}
  \item[拉伸 Extrude] 选中草图/线框/面，沿法向拉伸（默认 10 mm）；闭合轮廓自动成实体；可设反向长度、拔模角、对称。
        多选多个轮廓会各建一个。
  \item[旋转 Revolve] 轮廓绕轴旋转（默认 Z 轴过原点、360°）；选一条边即作为旋转轴。
  \item[放样 Loft] 至少两个截面，生成“实心”的过渡体；可设直纹、闭合。
  \item[扫掠 Sweep] 截面 + 路径：先选截面，最后选路径对象（或直接点它的边）；可开 Frenet、实心。
  \item[镜像 Mirror] 关于 XY/XZ/YZ 平面，或第二选择作为镜像面。
  \item[缩放 Scale] 默认均匀 1.0，可切换为 X/Y/Z 分向。
  \item[投影到表面] 先点目标面，再选要投影的形状/边。
\end{description}

\section{细节特征}

\begin{itemize}
  \item \ui{圆角}/\ui{倒角}：先问半径/尺寸（默认 1 mm），选中边则只处理所选边，未选则全部边；
        \textbf{如果选择集里有实体或实体内的特征，会自动改调 PartDesign 的圆角/倒角}（否则会生产一个挂在
        历史之外的形状）；
  \item \ui{厚度}（1 mm）、\ui{偏移}（3D/2D）、\ui{直纹面}、\ui{建立面}、\ui{剖面}；
  \item 树里右键 Part 圆角/倒角，可以 \ui{移入实体}，把它吸收进实体历史。
\end{itemize}

\section{阵列：LinkArray}

\ui{链接阵列} 一次提供五种：\textbf{线性、极坐标、圆周、路径、点}。第一个所选对象是被阵列的对象，
路径/点阵列取第二个所选对象作为参考。

\section{检查与测量}

\begin{itemize}
  \item \ui{检查几何}（\cmd{Part\_CheckGeometry}）与 \ui{形状信息} 共用一个信息面板：
        BOP 检查结果、实体/壳/面/线/边/点计数、体积/面积/长度、包围盒、质心；
  \item 测量不在零件工具条里：\menupath{工具 → 测量}——
        选一个元素给出坐标/边长/半径/面积/体积，选两个给出最小距离、ΔX/ΔY/ΔZ、两直线边或两个面的夹角；
  \item \ui{质量属性}：体积、面积、密度、质量、质心、惯性矩（质量需要先给对象赋材质密度）；
  \item \ui{单位换算}：单位之间的换算计算器。
\end{itemize}

\section{构建与复合}

\begin{itemize}
  \item \ui{形状构建器}（\cmd{Part\_Builder}）：两点成边、边成线框、点/边成面、面成壳、壳成实体；
  \item \ui{转为实体}：壳/复合体 → 实体；
  \item \ui{复合}、\ui{炸开复合}、\ui{复合过滤}、\ui{容差集}；
  \item \ui{去除特征}：点选面，直接去掉该面衍生的特征；
  \item \ui{精炼形状}、\ui{反转形状}、\ui{简单副本/变换副本/元素副本}；
  \item 选择过滤：顶点/边/面选择、\ui{选择过滤}、\ui{框选面}。
\end{itemize}

\section{导入与导出}

\ui{导入}/\ui{导出}走系统文件选择器，并按文件内容嗅探格式：

\begin{itemize}
  \item STEP / IGES → 零件导入；STL / PLY / OFF / 3MF → 网格；BREP → 零件；OBJ/IFC/DXF/SVG 走对应导入器；
  \item 导出时\textbf{改扩展名即改格式}：\file{.step}、\file{.iges}、\file{.brep}、\file{.stl}、\file{.obj}……；
  \item 导出多个对象前，可先 \ui{复合} 成一个对象；否则会分别写出。
\end{itemize}

\begin{practice}[练习：布尔求差与导出]
\begin{enumerate}
  \item 建一个立方体（20×20×20）和一个圆柱（半径 5、高 30）；
  \item 移动圆柱到立方体中心（树里双击圆柱改“放置”，或直接在对话框里改位置）；
  \item 先点立方体、再点圆柱 → \ui{布尔 → 剪切} → 得到带通孔的方块；
  \item \ui{检查几何} 确认 BOP 无误；
  \item \ui{文件 → 导出}，文件名写 \file{block.step}，保存。
\end{enumerate}
\end{practice}

\section{与零件设计的分工}

\begin{center}
\small
\begin{tabularx}{\linewidth}{@{}lX@{}}
\toprule
\textbf{场景} & \textbf{用哪个} \\
\midrule
要反复改尺寸、要历史 & 零件设计（PartDesign）\\
快速建毛坯、一次性运算 & 零件（Part）\\
处理导入的 STEP/IGES & 零件（Part）\\
把几个独立实体合成一个 & 零件（Part）的布尔/复合 \\
需要圆角且希望可回改 & 零件设计（PartDesign）\\
\bottomrule
\end{tabularx}
\end{center}

\chapter{曲面、网格、点云与其它工作台}

这一章覆盖“不是每天都用，但关键时刻能救场”的工作台：曲面补洞、网格修复、点云拟合、OpenSCAD、
逆向工程、检查与材质。它们都在 \ui{更多 ▾} 里。

\section{曲面（Surface）}

曲面工作台有 7 条命令，都是“先选边界，再建对象，再在对话框里调参数”：

\begin{center}
\small
\begin{tabularx}{\linewidth}{@{}lX@{}}
\toprule
\textbf{命令} & \textbf{作用} \\
\midrule
填充 Filling & 用一组边界曲线围出一块曲面；可调次数（默认 3）、曲线上点数（15）、迭代次数（2）\\
几何填充 GeomFillSurface & 在 2/3/4 条边界之间生成过渡曲面，可选填充类型 \\
截面 Sections & 多截面生成曲面（类似放样，但只出曲面）\\
延伸面 ExtendFace & 把已有面沿 U/V 方向延伸（默认 0.05 倍），可对称 \\
混合曲线 BlendCurve & 两条边之间的相切过渡曲线 \\
网格上的曲线 CurveOnMesh & 在网格上点选控制点，投影出样条曲线 \\
剪切 Cut & 两个形状互剪 \\
\bottomrule
\end{tabularx}
\end{center}

曲面对象不是实体：需要时用 \ui{零件 → 厚度} 加厚，或 \ui{转为实体}（缝合后）变成实体。

\begin{practice}[练习：给破洞的壳补面]
选洞口一圈边 → \ui{填充} → 调整次数使曲面平滑 → 确定 → 选中曲面与壳 → \ui{零件 → 转为实体}。
\end{practice}

\section{网格（Mesh）}

网格工作台处理三角网格：导入的 STL、3D 扫描结果、为分析而离散的模型。

\subsection{创建网格}

\begin{itemize}
  \item \ui{从形状建网格}（Standard）：线性偏差 0.1、角度 28.65°；
  \item \ui{Mefisto 网格}：按最大边长（默认 1.0）细分；
  \item \ui{从几何体建网格}：对形状做细分；
  \item \ui{规则实体}：立方体/圆柱/圆锥/球/椭球/圆环；
  \item \ui{Gmsh 重划网格}：使用内置 Gmsh（算法 + 最大/最小尺寸 + 面检测角）。
\end{itemize}

\subsection{修改与修复}

统一/翻转法线、补洞（按边数或点选）、加三角面（三点）、按组件数删除、平滑（Taubin/Laplace/中值，4 次）、
简化（减面率 0.9）、缩放、合并、按组件拆分。

\subsection{布尔与切割}

并/交/差（至少两个网格，用网格内核自身的布尔）；多边形切割/修剪/分段/按平面切割/剖面（先在视图上画
多边形，再选保留内或外）；按平面偏移修剪；截面（轴 + 条数）。

\subsection{分析}

\ui{网格分析} 汇总：点数、面数、是否封闭实体、非流形边、自交、翻转面，并提供一键修复按钮
（全部修复、去重点、修退化、去非流形、修朝向、修自交、去褶皱）。另有面评估、曲率、曲率信息、
实体评估、包围盒。

\begin{practice}[练习：STL 修复并转实体]
导入 STL → \ui{网格分析} → 看非流形/自交 → 点 \ui{全部修复} → \ui{平滑} 一次 →
\ui{零件 → 从网格建形状}（缝合容差 0.1）→ \ui{转为实体} → \ui{精炼形状} → 导出 STEP。
\end{practice}

\section{点云（Points）}

\begin{itemize}
  \item \ui{导入点云}：支持常见的点云文本/二进制格式；
  \item \ui{导出点云}：建议 \file{.asc}；
  \item \ui{转换为点云}：从形状或网格生成点云（点距 0.1）；
  \item \ui{结构化点云}：把点云整理成规则网格（可设列宽）；
  \item \ui{合并}：多个点云合并；
  \item \ui{多边形裁剪}：在视图上画多边形，保留外侧或内侧的点。
\end{itemize}

点云既是检查（Inspection）的“实测几何”，也是逆向工程拟合曲面的输入。

\section{OpenSCAD}

OpenSCAD 工作台提供 15 条命令，但\textbf{应用没有内置 OpenSCAD 可执行文件}，这一点必须知道：

\begin{itemize}
  \item \textbf{纯 Python 实现、可用}：镜像/缩放/调整网格特征、精炼形状、增加容差、边转面、展开放置、
        炸开组、替换对象、删除对象及其子对象、按颜色着色；
  \item \textbf{需要外部二进制、在移动端会失败}：添加 OpenSCAD 元素、凸包（Hull）、闵可夫斯基和、
        网格布尔。它们会报“OpenSCAD executable unavailable”。
  \item 第二条工具条直接复用零件命令：检查几何、图元、形状构建器、剪切/结合/交集、拉伸、旋转。
\end{itemize}

\section{逆向工程（Reverse Engineering）}

15 条命令中 13 条可用，核心是“从点/网格拟合出几何”：

\begin{itemize}
  \item \ui{近似平面/球/圆柱}：最小二乘拟合，输出对应曲面；
  \item \ui{近似曲面}：B 样条曲面，可设 U/V 次数与极点数；
  \item \ui{近似曲线}、\ui{近似多项式}（只输出系数与残差报告）；
  \item \ui{网格边界}：提取网格边界线框；
  \item \ui{分割}：按曲率容差（0.1）与每段最少面数（20）自动分块；
  \item \ui{最佳拟合分割}：按平面/圆柱/球拟合来分割；
  \item \ui{手动分割}：在网格上画多边形分块；
  \item 缺失：泊松重建（Poisson）与视图三角化（ViewTriangulation）。
\end{itemize}

\begin{practice}[练习：扫描件逆向]
导入点云 → \ui{裁剪} 去掉杂物 → \ui{近似球} 或 \ui{近似曲面} → 把结果加厚/转实体 → 导出。
\end{practice}

\section{检查（Inspection）}

用于曲面比对：

\begin{itemize}
  \item \ui{视觉检查}：多选“实测几何”和“名义几何”，设搜索距离（默认 0.05），生成检查组，
        在实测件上用蓝→绿→红画偏差色图；
  \item \ui{检查…}：在模型上点一下就读该点的距离；超出搜索半径会提示没有距离。
\end{itemize}

\section{材质（Material）}

\begin{itemize}
  \item \ui{编辑材质}：打开材质库，逐项修改物理属性与外观属性，“保存为新材质”进入用户库；
  \item \ui{设置材质}：给对象赋材质（同时按外观上色）；
  \item \ui{查看材质}/\ui{查看外观}：读取当前值；
  \item \ui{逐面上色}：给单个面上色；
  \item \ui{纹理映射}：给对象贴图。
\end{itemize}

\begin{tip}[质量属性需要密度]
\ui{零件 → 质量属性} 计算质量、质心、惯性矩时依赖材质的\textbf{密度}。先 \ui{设置材质}，再算质量。
\end{tip}

\chapter{绘图（Draft）：2D 线条与工作平面}

Draft 工作台适合画二维图：平面上的线、圆、多边形、标注，以及路径、阵列这类“图形式”的操作。
它的坐标系是\textbf{工作平面}，不是固定的 XY 平面——这是用好 Draft 的关键。

\section{入口与工具条}

从 \menupath{更多 ▾ → Draft} 进入，或使用 \ui{零件} 页里的 \ui{Draft} 工具组。工具按 FreeCAD 的
五条工具条分组：\ui{创建}、\ui{标注}、\ui{修改}、\ui{实用工具}、\ui{捕捉}。

\section{工作平面}

在 \ui{更多工具} 里选择 \ui{工作平面}，有五种：

\begin{center}
\small
\begin{tabularx}{\linewidth}{@{}lX@{}}
\toprule
\textbf{选项} & \textbf{含义} \\
\midrule
上 (XY) / 前 (XZ) / 侧 (YZ) & 三个坐标平面 \\
所选面 & 用当前选中的平面作为工作平面（先选中，否则提示“请先选择平面”；非平面会报错）\\
当前视图 & 平面过相机目标点，屏幕右/上为 u/v，视线为法向 \\
\bottomrule
\end{tabularx}
\end{center}

工作平面影响落点投影、矩形/圆/多边形的朝向、镜像与旋转轴、键入坐标的解释，以及栅格捕捉的取整方向。
它\textbf{不影响} 3D 网格的显示（网格始终在 XY 地面）。

\section{捕捉}

Draft 的捕捉条有 15 个开关，\ui{Lock} 是总开关：

\begin{itemize}
  \item \textbf{默认开启}：Lock、端点、中点、圆心、近点、交点、垂足、栅格、工作平面；
  \item \textbf{默认关闭}：角度、正交、平行、延伸、特殊、尺寸；
  \item 端点/中点/圆心取自可见形状：顶点、每条边的中点、圆/弧的中心，取 18 dp 内最近者；
  \item 没有对象可捕捉时，落点落在工作平面上：\ui{正交}锁定水平/竖直，否则 \ui{角度} 每 15° 吸附，
        最后 \ui{栅格} 按 1 mm 取整；
  \item 捕捉名称会显示在光标旁（例如“端点”“交点”）。
\end{itemize}

\section{点输入}

轻触落点；顶部横幅显示工具名、提示，以及 \ui{Typing}/\ui{撤销}/\ui{关闭}/\ui{完成}/\ui{复制}/✕。

\begin{itemize}
  \item \ui{Type}：输入 \file{x, y, z} 或 \file{@dx, dy, dz}（2 或 3 个数，逗号/分号/空格分隔）。
        坐标按\textbf{工作平面的局部坐标}解释，\file{@} 表示相对上一点；
  \item 长按进入精确模式（准星在手指上方 60 dp），抬起即落点；
  \item 多段工具双击结束（第二次轻触会撤掉刚加的点），定长工具画够就自动结束；
  \item \ui{完成} 需要达到最少点数；\ui{关闭} 只对折线且至少 2 点；\ui{撤销} 删掉最后一点。
\end{itemize}

\section{绘制工具}

\begin{longtable}{@{}p{0.18\linewidth}p{0.10\linewidth}p{0.62\linewidth}@{}}
\toprule \textbf{工具} & \textbf{点数} & \textbf{说明} \\ \midrule \endhead
线 & 2 & 起点 → 终点 \\
折线 & 多点 & 连点；双击结束 \\
矩形 & 2 & 对角两点 \\
圆 & 2 & 圆心 → 半径点 \\
圆弧 & 3 & 圆心 → 起点 → 终点 \\
三点圆弧 & 3 & 弧上三点 \\
椭圆 & 2 & 对角两点（长短轴自动交换）\\
多边形 & 2 & 先问边数（默认 6），再圆心 + 顶点（内接）\\
B 样条 & 多点 & 连点 \\
贝塞尔 & 多点 & 三次曲线，每段 4 点后继续 \\
点 & 1 & 单点 \\
文字 & 2 & 先输入文字（\file{\textbackslash n} 换行），再点位置 \\
形状字符串 & 2 & 再问高度（默认 10）；字体固定为 Roboto \\
标注 & 3 & 三点线性标注 \\
标签 & 多点 & 先输文字；目标取当前选择；再点目标点与标签点 \\
面 & — & 先选面再执行（面绑定）\\
填充 & — & 先选对象；选 8 种图案之一与比例（默认 10）\\
\bottomrule
\end{longtable}

\section{修改与阵列}

使用前\textbf{先选对象}，否则提示“请先选择对象”。

\begin{center}
\small
\begin{tabularx}{\linewidth}{@{}lX@{}}
\toprule
\textbf{工具} & \textbf{操作} \\
\midrule
移动 & 基点 → 目标点（横幅的“复制”开启则复制）\\
旋转 & 圆心 → 基准角点 → 旋转角点（轴为平面法向）\\
缩放 & 基点 → 输入因子（默认 2）\\
镜像 & 两点定镜像线 \\
偏移 & 点一下定侧与距离；圆/弧改半径，线按垂距；只处理第一个对象 \\
拉伸 & 框两点 → 基点 → 目标点 \\
修剪 / 延伸 & 点目标（面则沿法向拉伸，线则改最近端）\\
分割 & 先选线，再点分割位置 \\
编辑 & 点节点，再点新位置（仅对带 Points 的对象）\\
\bottomrule
\end{tabularx}
\end{center}

其它：\ui{升级}/\ui{降级}（会删原对象）、\ui{连接}、\ui{Draft↔Sketch 转换}、\ui{克隆}、
\ui{2D 形状视图}（固定 +Z）、\ui{线转 B 样条}、\ui{修复}、\ui{加入施工组}。

\subsection{阵列}

\begin{itemize}
  \item \ui{矩形阵列}：问 X/Y 份数与间隔（默认 2/2/100；Z 方向 1 且间隔与 X 相同）；
  \item \ui{极坐标阵列}：问元素数（5）与扫掠角（360°），\textbf{中心固定在原点}，不能点选；
  \item \ui{环形阵列}：问径向 100、切向 50、环数 3，中心同样在原点；
  \item \ui{路径阵列 / 扭曲路径阵列 / 点阵列}：按选择顺序（基体在前，路径或点对象在后），
        数量固定（4 或 15），没有参数对话框；
  \item Link 版阵列为链接阵列（引用原对象）。
\end{itemize}

\section{样式与图层}

\begin{itemize}
  \item \ui{设置样式}/\ui{应用当前样式}/\ui{标注样式编辑器} 都转到标准的“外观”对话框：
        \textbf{只有 12 色 + Default}，用于改对象颜色；移动端没有线宽、线型、图层样式与注释样式编辑器；
  \item 图层：\ui{新建图层}（有选择则把选择放入）、\ui{添加到图层}、\ui{图层管理器}（点图层切换显示）、
        \ui{自动分组}（设置新对象的默认图层）、\ui{添加到组}/\ui{选择组}；
  \item 施工模式：把新对象加入 Construction 组。
\end{itemize}

\begin{practice}[练习：带标注的平面布置图]
\begin{enumerate}
  \item 工作平面选 \ui{上 (XY)}，打开栅格捕捉；
  \item 用 \ui{矩形} 画 4000 × 3000 的房间（Type 输入对角坐标），再用 \ui{折线} 画一条隔墙；
  \item 用 \ui{标注} 标出两道主要尺寸；
  \item 用 \ui{矩形阵列} 复制一列柱网（X 间隔 1000、2 份）；
  \item 保存，转到 \ui{更多 → 零件} 用 \ui{拉伸} 把墙线变成 3D 墙体。
\end{enumerate}
\end{practice}

\section{源码可见的限制}

\begin{itemize}
  \item 样式只有颜色；极坐标/环形阵列的中心固定在原点；矩形阵列的 XYZ 共用一个间隔；
  \item 路径/点阵列没有参数对话框，且依赖选择顺序；
  \item 平行、延伸、特殊、尺寸这四个捕捉开关在源码中没有对应逻辑；
  \item 栅格间隔固定 1 mm，且网格画在 XY 地面而不是工作平面；
  \item 偏移只处理第一个对象；修剪/延伸只支持线/面；编辑只对有 Points 的对象；
  \item 形状字符串字体固定 Roboto；2D 形状视图固定 +Z；填充作用于整个对象。
\end{itemize}

\chapter{工程图（TechDraw）：从 3D 到二维蓝图}

工程图工作台把 3D 模型变成可交付的二维图纸：视图、剖视、局部放大、尺寸、符号，最后导出 PDF/SVG/DXF。
它是全流程（建模 → 装配 → 出图）的最后一站。

\section{页面与模板}

\begin{enumerate}
  \item 底部切到 \ui{工程图}；空状态显示 \ui{还没有工程图} 与 \ui{新建工程图页}；
  \item 点 \ui{新建工程图页}（命令 \cmd{TechDraw\_PageDefault}）；
  \item 页面使用偏好设置里的默认模板；没有设置时使用 A4 横向的 \file{Default\_Template\_A4\_Landscape.svg}；
  \item 模板字段（作者、日期、标题、比例、图号）可以被自动填充（取文档的 CreatedBy/Company/Label），
        也可以点字段直接改值。
\end{enumerate}

页面对象出现在特征树里；点页面行切换到它。双击视图打开属性（比例、方向、X/Y 位置等）；
\ui{重绘页面} 刷新所有视图。

\section{视图类型}

\begin{center}
\small
\begin{tabularx}{\linewidth}{@{}lX@{}}
\toprule
\textbf{视图} & \textbf{操作} \\
\midrule
视图 View & 选零件（可多选）→ 选方向（前/上/右/后/下/左/等轴测）→ 放在页面中心 \\
投影组 ProjectionGroup & 选零件 → 多选投影方向（默认前/上/右）→ 自动排成标准的第三角视图 \\
剖视图 SectionView & 先选基视图 → 选方向（右/左/上/下）→ 过视图中心剖切；符号自动 A/B/C，放在右（或下）侧 60 mm \\
复杂剖视 & 选基视图 → 选切割草图/线框 → 偏移策略 \\
局部详图 DetailView & 选基视图 → 在页面上点圆心 → 半径（默认 10）→ 自动生成 2× 放大、字母编号，放在 +80/+40 mm \\
活动视图 ActiveView & 当前 3D 视口的截图（宽 100 mm）\\
图像 Image & 插入图片（固定 100 × 100 mm）\\
符号 / 粗糙度 & 从资源 SVG 里选，插在页面中心 \\
注释 / 富文本 & 先输入文字，再点位置 \\
气球 Balloon & 选视图 → 点指向点 → 点气泡位置，编号自动 \\
引线 Leader & 选视图 → 点箭头点 → 继续点 → 完成 \\
焊接符号 & 选引线 → 选类型（当前版本类型参数未生效）\\
裁剪组 ClipGroup & 框一个区域并加/减视图（当前版本不实际裁剪）\\
\bottomrule
\end{tabularx}
\end{center}

还有：\ui{移动视图}、\ui{共享视图}（同一视图放多页）、\ui{叠放}、\ui{锁定}。

\section{尺寸标注}

在页面上点视图的边/顶点/面，再点尺寸命令。

\subsection{智能尺寸}

\ui{尺寸} 会按选择自动判断：

\begin{itemize}
  \item 选一个面 → 面积；
  \item 选两条边 → 角度；
  \item 选三个顶点 → 三点角；
  \item 选圆/圆弧边 → 询问直径还是半径；
  \item 其它 → 询问长度/水平/垂直。
\end{itemize}

\subsection{专用尺寸}

长度、水平、垂直、角度、三点角、面积、半径、直径、弧长、范围……每个命令都有选择校验，选错会提示
\ui{Incorrect selection: select …}。另有链式、坐标、倒角尺寸，以及基于所选尺寸的“对齐/层叠”。

尺寸会自动放置（一般偏 7 mm、半径 1.1r、角度 15 mm），\textbf{没有“点一下放尺寸”的步骤}；
放置后可直接拖动文字调整位置。

\subsection{格式}

增减小数位、加 ⌀/□/n× 前缀、去掉前缀、自定义 printf 格式、HoleShaftFit（H7/g6 等配合）、
\ui{修复尺寸}（把尺寸重新指向新的参考）。

\section{装饰}

\begin{itemize}
  \item 中心线（面/两条线/两个点）、修饰顶点/中点/象限点、修饰线/圆/圆弧、
        平行/垂直修饰线、橡皮擦、\ui{显示全部}；
  \item 线属性（线型 1--5、线宽 0.18--0.7）；
  \item 剖面线（Hatch，用 SVG 图案）与几何剖面线（Geometric Hatch，11 种 PAT 图案）。
\end{itemize}

\begin{warn}[装饰几何在页面上的显示有限]
当前版本里，通过“修饰”命令添加的中心线、修饰线等\textbf{不会绘制在页面上}；剖面线统一按斜线绘制，
线型/线宽也不生效。出图时请以视图轮廓 + 尺寸为主，装饰作为补充。
\end{warn}

\section{导出与打印}

\begin{itemize}
  \item \ui{导出为 SVG}：模板 + 轮廓 + 尺寸 + 文字 + 气球（引线、符号、图片、剖面线暂缺）；
  \item \ui{导出为 DXF}：二维 CAD 通用格式；
  \item \ui{导出为 PDF}：按页面 mm 换算成 72 dpi 的矢量绘制；入口在 \ui{更多工具} 里的
        \ui{将页面导出为 PDF}；
  \item \ui{Print All} 与 PDF 导出同一路径：\textbf{只导出当前页}，不经过系统打印对话框。
\end{itemize}

\section{页面交互}

\begin{itemize}
  \item \textbf{平移}：单指在页面空白处拖动；\textbf{缩放}：双指捏合（以两指中点为锚点）；右上角
        \ui{适合} 让整页适配屏幕；切页时自动适合；
  \item \textbf{移动}：拖动已选中的视图/尺寸/气球（视图改 X/Y，尺寸改文字位置，气球改气泡位置），
        一次拖动是一步可撤销操作；
  \item \textbf{选择优先级}：顶点 12 dp → 边 10 dp → 尺寸/气球锚点 18 dp → 面内部（最小面优先）→ 视图边框；
        要点选整个视图，点视图内空白处或在树里点；
  \item \textbf{高亮}：视图边深蓝加粗、顶点蓝点、面半透明蓝、视图虚线框；预选顶点青色；
  \item 提示语：“请先选择视图（或其边/顶点）”“请先在页面上创建页面”等。
\end{itemize}

\begin{practice}[练习：三视图 + 标注 + PDF]
\begin{enumerate}
  \item 打开第 4 章的支架（先保存）；
  \item 切 \ui{工程图} → \ui{新建工程图页}；
  \item \ui{投影组} → 选支架 → 投影选 \ui{前/上/右} → 生成三视图；
  \item 点视图边（变蓝加粗）→ \ui{长度}/\ui{水平}/\ui{垂直} 标外形尺寸；圆边 → \ui{直径}；
        两边 → \ui{角度}；拖尺寸文字调位置；
  \item \ui{更多工具 → 将页面导出为 PDF} → 选保存位置；
  \item 打开 PDF 检查：图纸、视图、尺寸都在。
\end{enumerate}
\end{practice}

\begin{practice}[练习：剖视与局部放大]
\begin{enumerate}
  \item 选前视图 → \ui{剖视图} → 方向“右” → 生成 A-A 剖视；
  \item 选中剖视，点内部面 → \ui{几何剖面线} → 选图案（如 Steel）；
  \item 选基视图 → \ui{局部详图} → 点放大中心 → 半径 10 → 得到 2× 的 Detail A；
  \item 在详图上加直径、中心线。
\end{enumerate}
\end{practice}

\section{源码可见的限制}

\begin{itemize}
  \item 打印只导出当前页，且没有系统打印对话框；
  \item SVG 导出缺引线、符号、图片、剖面线、修饰几何；
  \item 页面不绘制中心线/修饰几何/螺纹/装饰线，线型线宽不生效，剖面线统一斜线；
  \item 断裂视图没有断裂定义；裁剪组不裁剪；图像固定 100×100 mm；
  \item 剖视只能过视图中心四个方向、没有深度参数；局部详图固定 2× 圆形；
  \item 投影组没有第一角/第三角切换；视图没有“当前视图方向”或任意方向；
  \item 尺寸不能点放置；轴向长度退化为长度；弧长以线性尺寸表示；面积/弧长标注画成矩形气球；
  \item 页面比例/幅面只能通过模板或属性修改；“锁定”不阻止拖动。
\end{itemize}

\chapter{变量与表达式：让模型“会算”}

参数化建模的下一步，是把设计意图写进\textbf{变量}：把长、宽、孔径、阵列数集中到一张电子表格里，
所有尺寸引用它。改一处，整机更新。这就是“变量”选项卡（Spreadsheet 工作台）的作用。

\section{打开与创建}

\begin{enumerate}
  \item 切到底部的 \ui{变量} 选项卡；空状态显示 \ui{还没有电子表格} 与 \ui{创建电子表格}；
  \item 点 \ui{创建电子表格}（\cmd{Spreadsheet\_CreateSheet}）；
  \item 表格取代 3D 视图显示；树里出现 \ui{电子表格} 对象（默认名 \file{Spreadsheet}）。
\end{enumerate}

\section{单元格与输入}

\begin{itemize}
  \item \textbf{选中}：单击单元格；\textbf{区间}：长按锚定起点，再点对角单元格，得到 \file{A1:C4} 这样的区域；
  \item \textbf{编辑}：双击单元格，或用顶部公式栏右侧的 \ui{编辑} 按钮；
  \item 输入有三类：
    \begin{itemize}
      \item 文本：直接写；
      \item 带单位的数值：\file{25 mm}、\file{30 °}；
      \item 公式：以 \file{=} 开头，例如 \file{=A2*2}、\file{=Sheet2.B3}；
    \end{itemize}
  \item 公式栏左侧显示当前选区，右侧显示内容（单元格有别名时括注，例如“B3 (Thickness)”）。
\end{itemize}

\section{别名：给参数起名字}

\begin{enumerate}
  \item 选中一个单元格；
  \item 执行 \ui{设置别名}（\cmd{Spreadsheet\_SetAlias}）；
  \item 输入名字（例如 \file{Thickness}）。
\end{enumerate}

之后在\textbf{任何}数值属性里都可以写 \file{=Sheet.Thickness}（\file{Sheet} 是电子表格的对象名，
可在树里改）。如果别名在整个文档里唯一，也可以简写成 \file{=Thickness}。

\section{表达式：在哪里都能用}

属性行（包括特征对话框里的数值、草图尺寸、FEM 载荷、装配偏移……）都支持表达式：

\begin{itemize}
  \item 输入纯“数字 + 单位”：按数值写入；
  \item 输入其它任何内容（含运算、引用）：按\textbf{表达式}写入，属性行以蓝色 \file{=…} 显示；
  \item 尺寸与角度都支持；向量分量（例如放置的 X/Y/Z）同样走这套规则。
\end{itemize}

\begin{lstlisting}
=Sheet.Thickness            // 引用变量
=Sheet.Width - 2 * 3 mm     // 留边
=10 mm + 5 mm               // 直接运算
\end{lstlisting}

\section{格式与数据交换}

工具条还提供：六种对齐、加粗/斜体/下划线、合并/拆分单元格、列宽行高、插入/删除行列，
以及 \ui{导入}/\ui{导出} CSV/TSV（分隔符按扩展名取逗号、分号或制表符）。

\section{单位制}

\menupath{编辑 → 偏好设置 → 单位系统} 切换 FreeCAD 的单位方案与小数位。它影响\textbf{显示与输入解析}，
内部始终用毫米与千克：切换后已有数值不变，只是显示方式改变。

\begin{tip}[把参数表当“设计说明书”]
建议每个模型都建一张表：\file{Length}、\file{Width}、\file{Height}、\file{HoleDia}、\file{PatternCount}……
并在旁边写注释行说明含义。三个月后回来改设计，你会感谢自己。
\end{tip}

\begin{practice}[练习：把支架变成可配置的]
\begin{enumerate}
  \item 在变量页建表，A1 写 \file{Width}（值 60 mm）、A2 写 \file{Height}（40 mm）、A3 写 \file{Hole}（10 mm）；
  \item 选中 B1/B2/B3，分别设置别名 \file{Width}/\file{Height}/\file{Hole}；
  \item 回到支架草图，把两个尺寸改成 \file{=Sheet.Width}、\file{=Sheet.Height}；
  \item 把孔的直径改成 \file{=Sheet.Hole}；
  \item 回变量页把 B1 改成 100 mm——支架自动变宽，孔与圆角一起更新。
\end{enumerate}
\end{practice}

\section{限制}

\begin{itemize}
  \item 公式输入是弹窗 + 数字键盘，没有桌面版的函数补全下拉；
  \item 单元格类型（文本/数值/公式）按输入内容自动判定，没有显式类型设置；
  \item 大型表格没有筛选、排序、图表；它本质是“参数表”，不是 Excel。
\end{itemize}

\chapter{装配：把零件装到一起}

装配工作台把多个零件用\textbf{关节}（joint）连接起来，求解它们的相对位置，并检查自由度。它对应桌面
FreeCAD 的 Assembly 工作台，求解器是 OndselSolver。

\section{创建装配与插入零件}

\begin{enumerate}
  \item 切到 \ui{装配} 选项卡；
  \item 点 \ui{创建装配}（\cmd{Assembly\_CreateAssembly}）——生成装配容器与关节组；
  \item 点 \ui{插入零件}（\cmd{Assembly\_InsertLink}）：
        对话框列出\textbf{不在当前装配里}的实体（Body）、零件（App::Part）与形状（Part::Feature），
        选中即插入为链接（App::Link）；
  \item \textbf{第一个插入的零件会自动接地}（固定），这与桌面 FreeCAD 的行为一致。
\end{enumerate}

树里会看到 \ui{装配} → \ui{原点}、\ui{关节} 与各构件。

\section{关节}

\subsection{13 种关节类型}

固定（Fixed）、旋转（Revolute）、圆柱（Cylindrical）、滑动（Slider）、球（Ball）、距离（Distance）、
平行（Parallel）、垂直（Perpendicular）、角度（Angle）、齿轮齿条（RackPinion）、螺杆（Screw）、
齿轮（Gears）、皮带（Belt），另有 \ui{刚性组}（RigidGroup）。

\begin{tip}[没有单独的“平面”关节]
这个版本的引擎里，平面约束用“平行/垂直 + 距离”的组合来表达。
\end{tip}

\subsection{选择连接点}

\begin{enumerate}
  \item 点一个关节工具（例如 \ui{旋转关节}）；
  \item 在第一个零件上点\textbf{面、边或顶点}——必须点子元素，直接点零件会提示“请选择面、边或顶点”；
  \item 在\textbf{另一个零件}上点第二个连接点（点同一个零件会提示）；
  \item 两点齐全，立即创建关节并弹出关节对话框。
\end{enumerate}

连接器由“元素 + 最近顶点”生成（与桌面 TaskAssemblyCreateJoint 同一个算法）。对话框里可以：
重新选择 Reference1/2、改关节类型、改 \ui{偏移}（@Offset）、\ui{旋转}（@Rotation）、\ui{反向}（@Reverse）；
改完自动重新求解。

\section{求解与拖动}

\begin{itemize}
  \item \ui{求解装配}（\cmd{Assembly\_SolveAssembly}）：计算所有关节与自由度；失败时提示
        \ui{求解失败（代码）}；
  \item \textbf{拖动零件}：在视图中直接拖构件，应用先施加位移再求解——和桌面的拖拽器一样；
  \item \ui{接地}：选中零件点 \ui{切换固定}，再次执行取消。
\end{itemize}

\section{诊断}

装配出错时，用这些命令把问题“选中”出来：

\begin{itemize}
  \item \ui{选择冲突约束}、\ui{选择冗余约束}、\ui{选择格式错误约束}；
  \item \ui{有自由度的构件}、\ui{构件的关节}、\ui{选择链接对象}。
\end{itemize}

它们读取引擎记录的最后一次冲突/冗余/格式错误列表，把相关对象选中，便于在树里定位。

\section{输出与其它工具}

\begin{itemize}
  \item \ui{物料清单 BOM}（\cmd{Assembly\_CreateBom}）：生成物料清单对象（列可增删、拖动），可导出；
  \item \ui{快照}（\cmd{Assembly\_CreateSnapshot}）：冻结当前装配的副本；
  \item \ui{视图}（\cmd{Assembly\_CreateView}）：生成爆炸视图；
  \item \ui{仿真}（\cmd{Assembly\_CreateSimulation}）：运动仿真；
  \item \ui{导出 ASMT}：装配交换格式（走系统保存对话框）。
\end{itemize}

\begin{practice}[练习：做一个铰链]
\begin{enumerate}
  \item 建两个零件（底座与翻盖），都保存进同一个文档；
  \item 装配 → 创建装配 → 插入底座（自动接地）→ 插入翻盖；
  \item 在两个零件上各点一条同轴的边 → \ui{旋转关节}；
  \item \ui{求解装配}；拖动翻盖，确认它绕轴转动而不是乱跑；
  \item 如果拖动时其他方向也在动，检查是否漏了“重合/距离”约束，或添加平行/垂直约束补足自由度。
\end{enumerate}
\end{practice}

\section{限制}

\begin{itemize}
  \item 关节需要两次子元素选择，不能直接选整个零件；
  \item 没有独立的“平面”关节（用组合约束代替）；
  \item 大装配求解耗时随零件数增长，手机上建议分组件装配再总装。
\end{itemize}

\chapter{有限元（FEM）：在平板上做强度分析}

这是 MobileCAD 最“超出预期”的部分：它不是只能看结果的查看器——\textbf{CalculiX 求解器与 Gmsh 网格器
都打包在 APK 里}，在设备上完成建模、划网格、求解、看云图的全流程。

\section{设备上真正能跑什么}

\begin{center}
\small
\begin{tabularx}{\linewidth}{@{}lXX@{}}
\toprule
\textbf{组件} & \textbf{随 APK 打包} & \textbf{说明} \\
\midrule
CalculiX ccx 2.22 & 是（\file{libccx.so}） & 结构/热分析的求解器，真正在本机计算 \\
Gmsh 4.13.1 & 是（\file{libgmsh.so}） & 网格划分（体网格、局部细化）\\
OndselSolver & 是 & 装配求解器 \\
Netgen & 否 & 只能创建对象，无法用它划网格 \\
Elmer / Z88 / Mystran & 否 & 只有对象定义，没有可用的求解器 \\
VTK Python & 否 & 后处理改走内置的 FRD 解析与自绘管线 \\
\bottomrule
\end{tabularx}
\end{center}

\section{对象模型}

一次完整的分析由这些对象组成（都在特征树里）：

\begin{itemize}
  \item \ui{分析}（Analysis）容器：装材料、约束、网格、求解器、结果；
  \item \ui{材料}：固体材料（密度、杨氏模量、泊松比……）、流体材料、非线性材料等；
  \item \ui{边界条件与载荷}：固定、位移、力、压力、接触、弹簧、轴承、齿轮、温度、热流……
  \item \ui{网格}：Gmsh 网格、网格区域、Gmsh 细化；
  \item \ui{求解器}：\cmd{FEM\_SolverCalculiX}；
  \item \ui{结果}：CCX 结果对象与后处理管线。
\end{itemize}

\section{入门流程：十步完成一次分析}

\begin{enumerate}
  \item \textbf{建模}：在“零件”页做出实体（例如一根悬臂梁），保存；
  \item 切 \menupath{更多 ▾ → FEM}；
  \item \ui{分析}（\cmd{FEM\_Analysis}）：新建分析容器；
  \item 选中实体 → \ui{材料}（\cmd{FEM\_MaterialSolid}）→ 填密度、杨氏模量、泊松比；
  \item \ui{固定}（\cmd{FEM\_ConstraintFixed}）→ 选固定面；
  \item \ui{力}（\cmd{FEM\_ConstraintForce}）或 \ui{压力} → 选受力面 → 填大小与方向；
  \item \ui{从形状生成 Gmsh 网格}（\cmd{FEM\_MeshGmshFromShape}）→ 先给粗尺寸（如 5--10 mm）→ 应用；
  \item \ui{求解器}（\cmd{FEM\_SolverCalculiX}）→ \ui{运行求解器}（\cmd{FEM\_SolverRun}）；
  \item 选中结果对象 → \ui{显示结果}（\cmd{FEM\_ResultShow}）→ 选字段（等效应力 von Mises、
        位移、主应力、温度……）→ 设变形系数；
  \item 需要切片/曲线时，建 \ui{后处理管线} 并加过滤器（切面、等值面、沿线的数据、点数据、计算器）。
\end{enumerate}

\section{材料与边界条件}

\subsection{材料}

\ui{材料} 选固体材料后，在属性里填：密度、杨氏模量、泊松比、热膨胀系数、导热系数等。
手机上也可以从材质库选预置材料，再微调。

\subsection{边界条件与载荷}

工程上最常用的 16 种由应用直接创建（选中面/边自动加入 References）：固定、位移、力、压力、接触、
弹簧、轴承、齿轮、带轮、温度、热流、初始温度、刚体、变换、平面旋转、流体边界。
其余（Tie、自重、离心力、截面打印、电磁类、各类方程）走 FreeCAD 自带的 Python 命令。

\section{网格}

\begin{itemize}
  \item \ui{Gmsh 网格}：从形状生成，可设最大/最小单元尺寸、面检测角；
        任务面板的 \ui{应用} 就等于“生成网格”；
  \item \ui{网格区域}、\ui{网格组}、\ui{Gmsh 细化}：对局部加密；
  \item \ui{Mesh 重划 Gmsh}：对已有网格重新划分；
  \item \ui{清除网格}、\ui{网格信息}、\ui{网格转 Mesh}：管理与转换；
  \item \textbf{Netgen 不可用}——需要 Netgen 的命令在移动端会失败。
\end{itemize}

\begin{tip}[先粗后细]
手机 CPU 与散热有限。先用 5--10 mm 的粗网格验证边界条件是否正确，确认无误后再加密到 1--2 mm。
一次错误方向的细网格求解，可能白跑十几分钟。
\end{tip}

\section{求解}

\begin{itemize}
  \item 选中 \ui{求解器}（CalculiX）→ \ui{运行求解器}；
  \item 应用会写出 \file{.inp} 文件、调用设备上的 ccx、读取 \file{.frd} 结果；
  \item 求解期间会显示日志与进度；大模型请插电源、保持屏幕常亮；
  \item 失败时检查：单位是否一致、材料是否完整、约束是否足够（刚体位移）、网格是否有效。
\end{itemize}

\section{结果查看}

\ui{显示结果} 可选字段：等效应力（von Mises）、最大/最小主应力、位移、温度、剪应力……
并可设置\textbf{变形系数}（0 = 只看云图，1 = 真实变形，>1 = 放大）。

后处理管线与过滤器：\ui{从结果建管线}、\ui{变形过滤器}、\ui{等值面/裁剪}、\ui{切面}、
\ui{沿线数据}、\ui{点数据}、\ui{计算器}、\ui{分支过滤}，以及平面/球/圆柱/盒体等“创建函数”。

\screenshot[0.72]{tablet-fem.jpg}{有限元结果：树里是分析容器、材料、约束与 CCX 结果网格；视图里是应力云图。}

\begin{warn}[后处理的两个限制]
\ui{字形（Glyph）} 过滤器与 \ui{可视化} 命令在当前移植中不可用；结果是按表面顶点着色渲染的。
需要更复杂的可视化时，请导出结果到桌面 FreeCAD 处理。
\end{warn}

\begin{practice}[练习：悬臂梁]
\begin{enumerate}
  \item 建一个 100 × 10 × 10 mm 的长方体；
  \item FEM → 分析 → 材料（钢：密度 7850、E 210 GPa、ν 0.3）；
  \item 固定一端面；另一端面加 100 N 向下的力；
  \item Gmsh 网格尺寸 3 mm；
  \item 求解 → 看 von Mises 与位移；
  \item 把力改成 200 N，重新求解，观察应力翻倍。
\end{enumerate}
\end{practice}

\section{性能与限制小结}

\begin{itemize}
  \item 只支持 CalculiX 求解器（Elmer/Z88/Mystran 不可用）；Netgen 不可用；
  \item 无 Glyph/可视化过滤器；后处理以云图与基本管线为主；
  \item 手机内存有限：单元数建议控制在几十万以内，必要时用对称模型或壳单元；
  \item 求解器在设备上单线程/多线程运行，速度取决于设备散热。
\end{itemize}

\chapter{CAM 与 BIM：加工与建筑}

这两个工作台面向专业用途：CAM 生成数控刀路，BIM 做建筑信息模型。它们的命令都来自 FreeCAD 的官方
Python 模块，行为与桌面一致，但移动端各有取舍。

\section{CAM：从模型到 G-code}

\subsection{典型流程}

\begin{enumerate}
  \item 切 \menupath{更多 ▾ → CAM}；
  \item \ui{作业}（\cmd{CAM\_Job}）：用选中的实体（没选就用全部可见实体）建立加工任务，
        默认模板会建好作业坐标系与输出设置；
  \item 作业面板分三页：\ui{设置}、\ui{输出}、\ui{刀具}；
  \item 刀具：\ui{刀具控制器} → \ui{新建刀具}/\ui{从库加载}（刀具库随包提供）；
  \item 操作：\ui{轮廓}（Profile）、\ui{型腔}（Pocket Shape）、\ui{端面铣}（Mill Facing）、
        \ui{螺旋}（Helix）、\ui{自适应}（Adaptive）、\ui{槽}（Slot）、\ui{钻孔}（Drilling）、
        \ui{雕刻}（Engrave）、\ui{3D 型腔}（Pocket 3D）；
  \item 修饰：狗骨（Dogbone）、斜坡入刀（RampEntry）、标签（Tags）、引入引出（LeadInOut）、
        阵列/镜像、边界、Z 校正、轴映射、拖动刀、插铣等；
  \item \ui{校验}（Sanity）：检查刀路是否为空、进给/转速是否缺失；
  \item \ui{后处理}（\cmd{CAM\_Post}）：选后处理器（fanuc、grbl、linuxcnc、centroid、estlcam 等），
        生成 G-code 并通过系统文件选择器保存。
\end{enumerate}

\subsection{可用范围与限制}

\begin{itemize}
  \item 后处理器脚本随包提供，常见的三轴/四轴控制器都有；
  \item \textbf{没有材料去除仿真}：\cmd{CAM\_Simulator}/\cmd{CAMotics} 依赖的 CAMSimulator 未打包，
        只能在 3D 视图里查看刀路曲线，用 \ui{校验} 做基本检查；
  \item 复杂 3D 刀路（自适应、3D 型腔）在手机上计算较慢，建议先用粗公差预览。
\end{itemize}

\begin{practice}[练习：铣一个型腔]
建一个带圆角凹槽的方块 → CAM 作业 → 新建 φ6 平底刀 → 型腔操作选凹槽底面 → 设步距与下刀深度 →
校验 → 后处理（grbl）→ 保存 \file{pocket.nc}。
\end{practice}

\section{BIM：建筑与 IFC}

BIM 工作台复刻了 FreeCAD 的 BIM 模块（131 条命令），包括：

\begin{itemize}
  \item \textbf{建筑构件}：场地（Site）、建筑（Building）、楼层（Level）、空间（Space）、
        墙（Wall）、幕墙（CurtainWall）、柱（Column）、梁（Beam）、板（Slab）、门（Door）、
        窗（Window）、覆盖层（Covering）、管道（Pipe）、楼梯（Stairs）、屋顶（Roof）、
        面板（Panel）、框架（Frame）、围栏（Fence）、桁架（Truss）、设备（Equipment）、钢筋（Rebar）；
  \item \textbf{结构构件}：柱/梁/板、结构系统、钢筋、桁架、框架、幕墙；
  \item \textbf{材质}：材质、多材质、BIM 材质；
  \item \textbf{标注与图框}：尺寸、文字、轴网、剖面、2D 视图、图框；
  \item \textbf{管理}：项目设置、项目管理器、IFC 管理、图层、材质、预检。
\end{itemize}

\subsection{IFC 导入导出}

\begin{itemize}
  \item IfcOpenShell 0.8.4 随包提供；
  \item \textbf{打开 IFC}（\file{.ifc}）：以 NativeIFC 方式打开为工程文档；
  \item \textbf{导入 IFC}：作为对象插入当前文档（移动端默认“实体形状”策略、导入全部子项，
        因为手机上不能像桌面那样按需展开加载）；
  \item \textbf{导出 IFC}：\file{.ifc} 或 \file{.ifcJSON}；另有“创建 IFC 工程”“保存 IFC”“展开”
        “更新”“对比”等命令；
  \item 注意：IFC 使用完整 schema，文件与内存占用都较大，手机上建议先导入局部或简化模型。
\end{itemize}

\section{机器人（Robot）}

Robot 工作台用于运动学演示与轨迹规划：

\begin{itemize}
  \item 默认机器人是 KUKA IR500（内置 DH 参数；VRML 模型不在包里，按运动链自绘）；
  \item 可以建轨迹、路点（TCP 或点选）、从边生成轨迹、修饰、合成、归位、工具形状；
  \item \ui{仿真} 用逆解跑 60 帧；
  \item 可导出 KUKA \file{.src} 程序（走系统保存对话框）。
\end{itemize}

\section{其它专业工作台}

\ui{检查}（Inspection）与 \ui{材质}（Material）已在第 8 章介绍；\ui{逆向工程} 也在第 8 章。
这些工作台与 CAM/BIM 一样，命令覆盖完整，但受移动端屏幕与内存限制，适合做“查看、校核、小改动”，
大型项目仍建议回桌面完成。

\chapter{文件、格式与数据安全}

MobileCAD 没有自己的“云”，文件就是文件：通过 Android 的系统文件选择器（SAF）读写，位置由你决定。
本章讲清楚格式、保存流程、自动恢复与常见坑。

\section{原生文档：FCStd}

\begin{itemize}
  \item \file{.FCStd} 是 FreeCAD 的原生文档：保存\textbf{全部建模历史}（草图、特征、表达式、装配、图纸），
        体积小、可完全还原；
  \item 保存时会保留一份 \file{.fcstd1} 备份；
  \item 无头引擎默认不写 \file{GuiDocument.xml}（界面/显示信息），MobileCAD 会\textbf{保留并回写}
        颜色、透明度、面颜色、可见性——所以存回桌面打开，外观不会丢；
  \item 打开 \file{.FCStd} 时，若文档里有装配，应用会自动切到“装配”页。
\end{itemize}

\section{导入与导出}

\subsection{导入}

\ui{文件 → 导入} 或 \ui{零件 → 导入}，按内容嗅探格式：

\begin{center}
\small
\begin{tabularx}{\linewidth}{@{}lX@{}}
\toprule
\textbf{格式} & \textbf{进入方式} \\
\midrule
STEP / IGES & 零件导入（Import）\\
STL / PLY / OFF / AMF / 3MF & 网格（Mesh）\\
BREP & 零件（Part）\\
IFC & NativeIFC（BIM）\\
DXF / SVG / OBJ / DAE / CSG / 3DS / SHP 等 & FreeCAD 注册的导入器 \\
FEM 结果（\file{.frd}、Z88、Fenics） & 有限元结果导入 \\
\bottomrule
\end{tabularx}
\end{center}

\subsection{导出}

\ui{文件 → 导出}：\textbf{改扩展名即改格式}——

\begin{itemize}
  \item 实体：STEP（\file{.step/.stp}）、IGES、BREP、STL、OBJ、PLY、OFF、AMF、3MF；
  \item 建筑：IFC、IFCJSON；
  \item 二维/图纸：DXF、SVG（扁平化）、PDF；
  \item 其它：WebGL（\file{.html}）、CSG/SCAD、JSON；
  \item 图像：截图 PNG；工程图页导出 PDF/SVG/DXF；电子表格 CSV/TSV；CAM G-code；机器人 \file{.src}。
\end{itemize}

\begin{warn}[两件“不能”的事]
\begin{itemize}
  \item \textbf{DWG 不能打开}：FreeCAD 的 DWG 导入依赖外部的 ODA/LibreDWG/QCAD 转换器，APK 没有打包；
        请先转成 DXF。
  \item \textbf{glTF 导入可能不可用}：运行库缺少对应的导入模块；需要时先转成 STEP/STL。
\end{itemize}
\end{warn}

\section{系统文件选择器（SAF）}

所有文件操作都走 Android 的 SAF：

\begin{enumerate}
  \item 你触发“打开/保存/导入/导出”后，应用先把请求记下来（\ui{FileRequest}），再拉起系统选择器；
  \item \textbf{打开}：把选中的文件复制到应用缓存，再交给引擎；应用会申请持久读权限，让“最近文件”在重启后仍可用；
  \item \textbf{保存}：引擎先写缓存文件，再通过输出流写回你选的 URI；
  \item 文件选择器不过滤扩展名（只显示 \file{*/*}），所以\textbf{保存时务必自己写扩展名}；
  \item 打开时没有扩展名（例如内容提供者给出 \file{document:1234}），应用会按内容嗅探：
        STEP 头、IGES、BREP、STL 长度校验、FCStd 的 \file{Document.xml}、3MF、SVG、DXF、OBJ。
\end{enumerate}

\begin{tip}[保存的两种方式]
\begin{itemize}
  \item \menupath{文件 → 保存}：已知位置时\textbf{静默覆盖}，不弹选择器；
  \item \menupath{文件 → 另存为}：总是弹选择器，指定新位置/新名字。
\end{itemize}
\end{tip}

\section{最近文件与其它文档命令}

\begin{itemize}
  \item \ui{最近文件}：最多保留 \textbf{8} 条，存在应用配置里；可从这里快速重开；
  \item \ui{另存副本}（SaveCopy）：导出当前文档的副本（同样保留视图数据）；
  \item \ui{合并项目}（MergeProjects）：把另一个 \file{.FCStd} 的对象合并进当前文档；
  \item \ui{文档信息}（ProjectInfo）：查看/修改标题、作者、公司、许可、注释等（其中 6 项可改）；
  \item \ui{恢复到已保存版本}（Revert）：丢弃未保存修改，重新载入磁盘版本。
\end{itemize}

\section{自动保存与崩溃恢复}

\begin{itemize}
  \item 每 \textbf{60 秒}、以及应用切到后台时，若文档有改动，会把恢复副本写到应用私有目录
        \file{files/recovery/}（\file{recovery.FCStd} + 记录 URI/名称/时间的 JSON）；
  \item 正常保存后恢复副本被删除；
  \item 应用异常退出后再次启动，会自动打开恢复副本并提示
        \ui{已恢复未保存的工作 from HH:mm}；此时请立即“保存”到正式文件。
\end{itemize}

\section{单文档与内存}

MobileCAD 是\textbf{单文档界面}：新开文件会替换当前视图（旧文档仍在内存，可能被系统回收）。大项目建议：
\begin{itemize}
  \item 一个项目一个 \file{.FCStd}，必要时拆成多个文件用“合并项目”汇总；
  \item 频繁保存，重要节点“另存副本”；
  \item 切换应用前先保存。
\end{itemize}

\section{单位制}

所有长度内部以毫米、质量以千克存储。\menupath{偏好设置 → 单位系统} 只改变显示与输入解析（例如换成英寸），
不会改变模型的实际尺寸。

\chapter{效率技巧：像老手一样操作}

这一章收集的是“知道与不知道差别很大”的操作习惯。每条都不长，但能让你的建模速度翻倍。

\section{手势的效率用法}

\begin{center}
\small
\begin{tabularx}{\linewidth}{@{}lX@{}}
\toprule
\textbf{场景} & \textbf{技巧} \\
\midrule
手指挡住目标 & 长按进入\textbf{精密指针}：准星出现在手指上方 60 dp，抬起手指即完成拾取 \\
选不准小元素 & 先用两指放大，再轻触；拾取有优先级（顶点 12 dp → 边 10 dp → 面） \\
需要旋转草图视角 & 草图中单指用于绘制，\textbf{三指}才能旋转；或者点导航立方体 \\
误触清空选择 & 双击是“清除选择”，做精确选择时少用双击 \\
想右键 & \textbf{双指轻点}视图，弹出与桌面右键等价的菜单 \\
批量选对象 & \textbf{双指长按}进入框选，拖动矩形；配合选择过滤器只框面或只框对象 \\
退出当前工具 & 右下角的工具条上有 ✕（等同 Esc）；或按返回键 \\
\bottomrule
\end{tabularx}
\end{center}

\section{返回键 = 万能 Esc}

记住这条就够：\textbf{返回键按层关闭}，从最内层开始——数字键盘、文本输入、选择器、菜单、工具面板、Draft、
装配配合、任务面板、特征对话框、草图。全部关闭后应用退到后台，不会退出。所以“卡在某个模式里”时，
连按返回键总能回到正常状态。

\section{选择与选择过滤器}

\begin{itemize}
  \item 轻触是\textbf{累加选择}，再点同一元素取消；选择条上的 ✕ 移除单项；双击清空。
  \item 需要只选面或只选边：\menupath{工具} 里的选择过滤（或命令 \cmd{Part\_FaceSelection} 等），
        视图右下角会出现“仅选择面”的提示条，点 ✗ 清除过滤。
  \item 选择历史：\menupath{编辑 → 选择后退/前进}（\cmd{Std\_SelBack}/\cmd{Std\_SelForward}）可以在选择集之间来回。
\end{itemize}

\section{用表达式代替数字}

任何数值输入都能写表达式：数字键盘最上面直接输入，或以 \file{=} 开头。例如：

\begin{lstlisting}
=60 mm / 2          // 30 mm
=Sheet.Thickness*2  // 引用电子表格里的别名（第 11 章）
=10 mm + 5 mm       // 带单位的运算
\end{lstlisting}

好处是：改一个变量，所有引用它的尺寸一起变。复杂模型里，把关键尺寸集中放到电子表格（变量选项卡），
再用表达式引用，是最省心的做法。

\section{用 Tip 回放建模历史}

实体（Body）的 \ui{Tip} 标记决定“算到哪一步”。在树里长按某个特征，选 \ui{将 Tip 滚动到此处}，模型会回退到
那一步的状态：可以检查中间结果、在早期特征上做修改。再滚回最后的特征即可恢复完整模型。

配合 \ui{压缩}（Suppress）可以临时排除某个特征，对比“有它/没它”的结果；排除带来的错误也会一并消失。

\section{工具搜索与最近工具}

\begin{itemize}
  \item 工具面板顶部的 \ui{搜索所有工具} 是全局搜索：输入中文命令名（凸台）、英文名（pad）或命令 ID
        （PartDesign\_Pad）都能命中；
  \item 工具条右侧的“最近工具”保留 6 个，反复用的命令会留在手边；
  \item 长按工具面板里的按钮不会弹出说明，命令一词即是 FreeCAD 的命令名，可在附录 A 与网上文档对照。
\end{itemize}

\section{把常用面板放在“变量”选项卡}

电子表格不只是算数：把设计参数（长、宽、高、孔径、阵列数）写成一张表，建模时全部引用它。改一处，
整机更新。这是参数化设计的“驾驶舱”。

\section{性能与流畅度}

\begin{itemize}
  \item \textbf{隐藏不看的对象}：大装配里把无关构件隐藏（树里长按 → 隐藏），视图重绘会明显变快；
  \item \textbf{减少同时打开的对话框}：特征对话框实时预览，复杂模型上会持续计算；
  \item \textbf{分部件建模}：几百个特征的单一实体不如拆成几个实体再装配；
  \item \textbf{有限元网格}：移动端建议先用较粗网格试算，确认边界条件后再加密；求解时插上电源并让屏幕常亮；
  \item \textbf{存储}：运行时约 1 GB 常驻，剩余空间不足 2 GB 时系统可能杀掉后台应用。
\end{itemize}

\section{硬件键盘与鼠标}

平板接上键盘/鼠标后，输入框可直接键入，布局不重建（插入键盘不会打断当前操作）。鼠标左键=轻触，
右键=上下文菜单，滚轮=缩放（等价于双指捏合）。触控板的多指手势由系统处理。

\section{自动保存与恢复}

MobileCAD 每 60 秒（以及每次切到后台时）把当前文档写入恢复副本。若应用异常退出，下次启动会自动打开
恢复副本并提示“已恢复未保存的工作 from HH:mm”；这时请立刻“保存”到一个正式文件。正常保存后恢复副本会被删除。

\chapter{故障排除与数据安全}

\section{启动问题}

\begin{description}
  \item[启动画面卡在“正在解压 CAD 引擎…”] 首次启动或应用升级后需要解压约 1 GB 运行时。请保持屏幕常亮、
        确保剩余空间 ≥ 2 GB。若超过 3 分钟仍无进展，强制停止应用后重开；仍失败则卸载重装（文档在存储卡上不受影响）。
  \item[红字“CAD 引擎启动失败：…”] 记下红字后面的错误信息。常见原因：存储空间不足、设备 ABI 不匹配
        （本 APK 只支持 arm64）、系统 WebView/运行库异常。可先重启设备再试；持续失败请带错误信息反馈。
  \item[启动后立即闪退] 检查是否安装了与官方版签名冲突的重打包版本；卸载后重装对应渠道的包。
\end{description}

\section{特征失败}

特征树里出现红色文字/叹号，说明该特征无法计算。处理顺序：

\begin{enumerate}
  \item 双击该特征打开对话框，看顶部红条的错误文字（例如“无法在基准特征上找到该几何”“基体特征上不存在该几何”）；
  \item 检查它的\textbf{引用}：轮廓是否闭合、是否自相交、引用的面/边是否还存在（前面的特征改过形，可能让引用失效）；
  \item 用 \ui{压缩} 临时跳过它，确认后面的特征是否依赖它；
  \item 回到草图用 \ui{约束} 列表排查（冲突/冗余/格式错误），或删掉多余线段；
  \item 修复后，长按该特征选 \ui{重新计算对象}。
\end{enumerate}

\section{草图问题}

\begin{description}
  \item[状态条一直显示“欠约束 N 个自由度”] 还差 N 个约束。常见遗漏：没有把某个点固定在原点附近、缺少水平/竖直约束、
        缺少一个尺寸。用 \ui{约束} 列表逐条查看。
  \item[冲突（Conflicting）] 两条约束互相矛盾（例如同时给了水平与竖直距离又给了角度）。删除其中一条。
        冲突约束会显示红色，\ui{约束} 面板可只看“冲突/冗余”。
  \item[冗余（Redundant）] 约束重复（同一条边被约束两次），删除多余的那条。
  \item[草图跳形] 欠约束时拖动几何或改尺寸，形状可能整体翻转。把草图约束到“完全约束”并给关键点加锁定。
  \item[尺寸是驱动还是参考] 尺寸可切换为“参考（从动）”：它只显示测量值、不参与约束，用于标注检查。
\end{description}

\section{文件与导出}

\begin{description}
  \item[打开的文件是空白/报错] 确认文件格式受支持；\file{.FCStd} 建议用相近版本的 FreeCAD 保存。
        从网上下载的 STEP/IGES 若结构复杂，导入可能耗时。
  \item[导出 STEP 后没有实体] 模型里可能只有网格或曲面；先用“零件 → 检查几何”确认是有效实体。
  \item[STL 表面很粗糙] 降低导出时的偏差（0.01--0.05 mm）；或先加密网格。
  \item[导出 PDF 是空白页] 工程图页面里还没有视图，或视图比例过大跑出纸面；检查页面与视图属性。
\end{description}

\section{崩溃与恢复}

MobileCAD 有崩溃恢复：每 60 秒自动保存恢复副本。异常退出后重新打开，应用会提示恢复。建议：
\begin{itemize}
  \item 重要节点手动 \menupath{文件 → 保存}；
  \item 用 \menupath{文件 → 另存副本} 做版本快照；
  \item 如果恢复提示没有出现，检查应用私有目录下的 \file{files/recovery/}（需 root 或调试版）。
\end{itemize}

\section{更新问题}

\begin{description}
  \item[提示“更新失败：校验和不匹配”] 下载损坏，重试即可；应用会校验 SHA-256。
  \item[提示“更新未安装”] 需要允许 MobileCAD 安装未知应用：按提示进入系统设置授权后返回。
  \item[F-Droid 版不提示更新] 正常：F-Droid 版本通过 F-Droid 更新，应用内会提示“更新来自安装来源”。
\end{description}

\section{性能与内存}

\begin{itemize}
  \item 大装配先隐藏无关对象；复杂模型避免同时打开多个实时预览对话框；
  \item 有限元求解会占满 CPU，注意设备发热与降频；网格数先从几千开始；
  \item 剩余存储低于 2 GB 时系统可能回收后台，请先保存再切换应用。
\end{itemize}

\section{反馈与求助}

\begin{itemize}
  \item 项目仓库：\url{https://github.com/LucasCoraca/mobilecad}（issue 里请附：设备型号、Android 版本、
        MobileCAD 版本、FreeCAD 提交号、复现步骤、截图或录屏）；
  \item FreeCAD 本身的问题（几何内核、算法）同样可以反馈给上游 FreeCAD 项目；
  \item 本手册的勘误：在仓库提 issue 并标注 “manual”。
\end{itemize}


\appendix
\chapter{常用命令速查}
\markboth{附录 A 常用命令速查}{附录 A 常用命令速查}

下表按工作台列出主要命令（中文、英文、命令 ID），优先收工具条上可见的命令。命令 ID 与 FreeCAD 桌面
完全一致：在桌面版里搜索同一个 ID，或在 Python 控制台执行 \file{Gui.runCommand("ID")}，效果相同。
应用内工具面板的\ui{搜索所有工具}可检索全部 984 条命令，数据源为仓库中的 \file{assets/freecad/catalog.json}。

\section{核心（Std）}
\begin{longtable}{@{}p{0.32\linewidth}p{0.30\linewidth}p{0.30\linewidth}@{}}
\toprule 中文 & English & 命令 ID \\ \midrule \endhead
关于 \%1 & About \%1 & \cmd{Std\_About} \\
关于 Qt & About Qt & \cmd{Std\_AboutQt} \\
激活窗口 & Activate Window & \cmd{Std\_WindowsMenu} \\
对齐到所选 & Align to Selection & \cmd{Std\_AlignToSelection} \\
对齐到… & Align To... & \cmd{Std\_Alignment} \\
注释标签 & Annotation Label & \cmd{Std\_AnnotationLabel} \\
外观 & Appearance & \cmd{Std\_SetAppearance} \\
附加到远程调试器 & Attach to Remote Debugger & \cmd{Std\_MacroAttachDebugger} \\
底视 & Bottom & \cmd{Std\_ViewBottom} \\
包围盒 & Bounding Box & \cmd{Std\_SelBoundingBox} \\
框选元素 & Box Element Selection & \cmd{Std\_BoxElementSelection} \\
框选 & Box Selection & \cmd{Std\_BoxSelection} \\
框选缩放 & Box Zoom & \cmd{Std\_ViewBoxZoom} \\
可分割条 & Breakable Bar & \cmd{Std\_TestProgress1} \\
覆盖面板穿透鼠标事件 & Bypass Mouse Events in Overlay Panels & \cmd{Std\_DockOverlayMouseTransparent} \\
层叠 & Cascade & \cmd{Std\_CascadeWindows} \\
选择打开的窗口 & Choose Open Window & \cmd{Std\_Windows} \\
澄清选择 & Clarify Selection & \cmd{Std\_ClarifySelection} \\
剖切视图 & Clipping View & \cmd{Std\_ToggleClipPlane} \\
关闭 & Close & \cmd{Std\_CloseActiveWindow} \\
全部关闭 & Close All & \cmd{Std\_CloseAllWindows} \\
折叠选定项目 & Collapse Selected Items & \cmd{Std\_TreeCollapse} \\
折叠/展开 & Collapse/Expand & \cmd{Std\_TreeCollapseDocument} \\
命令行 & Command Line & \cmd{Std\_CommandLine} \\
复制 & Copy & \cmd{Std\_Copy} \\
自定义 & Customize & \cmd{Std\_DlgCustomize} \\
剪切 & Cut & \cmd{Std\_Cut} \\
删除 & Delete & \cmd{Std\_Delete} \\
依赖关系图 & Dependency Graph & \cmd{Std\_DependencyGraph} \\
开发人员手册 & Developers Handbook & \cmd{Std\_DevHandbook} \\
\bottomrule
\end{longtable}
\section{草图（Sketcher）}
\begin{longtable}{@{}p{0.32\linewidth}p{0.30\linewidth}p{0.30\linewidth}@{}}
\toprule 中文 & English & 命令 ID \\ \midrule \endhead
圆弧 & Arc & \cmd{Sketcher\_CompCreateArc} \\
贝塞尔曲线 & B-Spline & \cmd{Sketcher\_CompCreateBSpline} \\
圆锥 & Conic & \cmd{Sketcher\_CompCreateConic} \\
尺寸 & Dimension & \cmd{Sketcher\_CompDimensionTools} \\
编辑边 & Edit Edges & \cmd{Sketcher\_CompCurveEdition} \\
外部几何元素 & External Geometry & \cmd{Sketcher\_CompExternal} \\
圆角/倒角 & Fillet/Chamfer & \cmd{Sketcher\_CompCreateFillets} \\
水平/垂直约束 & Horizontal/Vertical Constraint & \cmd{Sketcher\_CompHorVer} \\
多边形 & Polygon & \cmd{Sketcher\_CompCreateRegularPolygon} \\
折线 & Polyline & \cmd{Sketcher\_CompLine} \\
半径/直径尺寸 & Radius/Diameter Dimension & \cmd{Sketcher\_CompConstrainRadDia} \\
矩形 & Rectangle & \cmd{Sketcher\_CompCreateRectangles} \\
长圆槽 & Slot & \cmd{Sketcher\_CompSlot} \\
切换约束 & Toggle Constraints & \cmd{Sketcher\_CompToggleConstraints} \\
将视图对齐到草图 & Align View to Sketch & \cmd{Sketcher\_ViewSketch} \\
角度尺寸 & Angle Dimension & \cmd{Sketcher\_ConstrainAngle} \\
三点创建圆弧 & Arc From 3 Points & \cmd{Sketcher\_Create3PointArc} \\
从圆心创建圆弧 & Arc From Center & \cmd{Sketcher\_CreateArc} \\
弧形槽 & Arc Slot & \cmd{Sketcher\_CreateArcSlot} \\
附加草图 & Attach Sketch & \cmd{Sketcher\_MapSketch} \\
通过节点创建B样条 & B-Spline From Knots & \cmd{Sketcher\_CreateBSplineByInterpolation} \\
固定约束 & Block Constraint & \cmd{Sketcher\_ConstrainBlock} \\
取消编辑 & Cancel Editing & \cmd{Sketcher\_CancelSketch} \\
影子克隆 & Carbon Copy & \cmd{Sketcher\_CarbonCopy} \\
居中矩形 & Centered Rectangle & \cmd{Sketcher\_CreateRectangle\_Center} \\
倒角 & Chamfer & \cmd{Sketcher\_CreateChamfer} \\
三点创建圆 & Circle From 3 Points & \cmd{Sketcher\_Create3PointCircle} \\
从圆心创建圆 & Circle From Center & \cmd{Sketcher\_CreateCircle} \\
克隆 & Clone & \cmd{Sketcher\_CompCopy} \\
重合约束 & Coincident Constraint & \cmd{Sketcher\_ConstrainCoincidentUnified} \\
\bottomrule
\end{longtable}
\section{零件设计（PartDesign）}
\begin{longtable}{@{}p{0.32\linewidth}p{0.30\linewidth}p{0.30\linewidth}@{}}
\toprule 中文 & English & 命令 ID \\ \midrule \endhead
增料图元 & Additive Primitive & \cmd{PartDesign\_CompPrimitiveAdditive} \\
创建基准 & Create Datum & \cmd{PartDesign\_CompSketches} \\
减料图元 & Subtractive Primitive & \cmd{PartDesign\_CompPrimitiveSubtractive} \\
增料立方体 & Additive Box & \cmd{PartDesign\_AdditiveBox} \\
增料圆锥体 & Additive Cone & \cmd{PartDesign\_AdditiveCone} \\
增料圆柱体 & Additive Cylinder & \cmd{PartDesign\_AdditiveCylinder} \\
增料椭球体 & Additive Ellipsoid & \cmd{PartDesign\_AdditiveEllipsoid} \\
添加式螺旋 & Additive Helix & \cmd{PartDesign\_AdditiveHelix} \\
添加式放样 & Additive Loft & \cmd{PartDesign\_AdditiveLoft} \\
添加式管道 & Additive Pipe & \cmd{PartDesign\_AdditivePipe} \\
增料棱柱体 & Additive Prism & \cmd{PartDesign\_AdditivePrism} \\
增料球体 & Additive Sphere & \cmd{PartDesign\_AdditiveSphere} \\
增料圆环体 & Additive Torus & \cmd{PartDesign\_AdditiveTorus} \\
增料楔形体 & Additive Wedge & \cmd{PartDesign\_AdditiveWedge} \\
布尔运算 & Boolean Operation & \cmd{PartDesign\_Boolean} \\
倒角 & Chamfer & \cmd{PartDesign\_Chamfer} \\
圆周阵列 & Circular Pattern & \cmd{PartDesign\_CircularPattern} \\
克隆 & Clone & \cmd{PartDesign\_Clone} \\
基准线 & Datum Line & \cmd{PartDesign\_Line} \\
基准面 & Datum Plane & \cmd{PartDesign\_Plane} \\
基准点 & Datum Point & \cmd{PartDesign\_Point} \\
清除特征 & Defeaturing & \cmd{PartDesign\_Defeaturing} \\
Draft & Draft & \cmd{PartDesign\_Draft} \\
复制对象 & Duplicate Object & \cmd{PartDesign\_DuplicateSelection} \\
圆角 & Fillet & \cmd{PartDesign\_Fillet} \\
挖槽 & Groove & \cmd{PartDesign\_Groove} \\
孔 & Hole & \cmd{PartDesign\_Hole} \\
渐开线齿轮 & Involute Gear & \cmd{PartDesign\_InvoluteGear} \\
线性阵列 & Linear Pattern & \cmd{PartDesign\_LinearPattern} \\
局部坐标系 & Local Coordinate System & \cmd{PartDesign\_CoordinateSystem} \\
\bottomrule
\end{longtable}
\section{零件（Part）}
\begin{longtable}{@{}p{0.32\linewidth}p{0.30\linewidth}p{0.30\linewidth}@{}}
\toprule 中文 & English & 命令 ID \\ \midrule \endhead
复合工具 & Compound Tools & \cmd{Part\_CompCompoundTools} \\
合并形状 & Join Shapes & \cmd{Part\_CompJoinFeatures} \\
链接阵列 & Link Arrays & \cmd{Part\_LinkArrays} \\
偏移 & Offset & \cmd{Part\_CompOffset} \\
分割形状 & Split Shapes & \cmd{Part\_CompSplitFeatures} \\
2D 偏移 & 2D Offset & \cmd{Part\_Offset2D} \\
3D 偏移 & 3D Offset & \cmd{Part\_Offset} \\
每面外观 & Appearance per Face & \cmd{Part\_ColorPerFace} \\
附着 & Attachment & \cmd{Part\_EditAttachment} \\
布尔碎片 & Boolean Fragments & \cmd{Part\_BooleanFragments} \\
布尔运算 & Boolean Operation & \cmd{Part\_Boolean} \\
布尔异或 & Boolean XOR & \cmd{Part\_XOR} \\
盒子修复1 & Box Fix 1 & \cmd{Part\_Box2} \\
盒子修复2 & Box Fix 2 & \cmd{Part\_Box3} \\
框选 & Box Selection & \cmd{Part\_BoxSelection} \\
倒角 & Chamfer & \cmd{Part\_Chamfer} \\
检查几何 & Check Geometry & \cmd{Part\_CheckGeometry} \\
圆周链接阵列 & Circular Link Array & \cmd{Part\_LinkArrayCircular} \\
组合 & Compound & \cmd{Part\_Compound} \\
组合过滤器 & Compound Filter & \cmd{Part\_CompoundFilter} \\
圆锥体 & Cone & \cmd{Part\_Cone} \\
连接形状 & Connect Shapes & \cmd{Part\_JoinConnect} \\
转换为实体 & Convert to Solid & \cmd{Part\_MakeSolid} \\
坐标系 & Coordinate System & \cmd{Part\_CoordinateSystem} \\
横截面 & Cross-Sections & \cmd{Part\_CrossSections} \\
立方体 & Cube & \cmd{Part\_Box} \\
剪切 & Cut & \cmd{Part\_Cut} \\
切出形状 & Cutout Shape & \cmd{Part\_JoinCutout} \\
圆柱体 & Cylinder & \cmd{Part\_Cylinder} \\
基准线 & Datum Line & \cmd{Part\_DatumLine} \\
\bottomrule
\end{longtable}
\section{绘图（Draft）}
\begin{longtable}{@{}p{0.32\linewidth}p{0.30\linewidth}p{0.30\linewidth}@{}}
\toprule 中文 & English & 命令 ID \\ \midrule \endhead
二维视图工具 & 2D View Tools & \cmd{Draft\_Shape2DViewTools} \\
圆弧工具 & Arc Tools & \cmd{Draft\_ArcTools} \\
阵列工具 & Array Tools & \cmd{Draft\_ArrayTools} \\
贝塞尔工具 & Bézier Tools & \cmd{Draft\_BezierTools} \\
添加到构造组 & Add to Construction Group & \cmd{Draft\_AddConstruction} \\
添加到组 & Add to Group & \cmd{Draft\_AddToGroup} \\
添加到图层 & Add to Layer & \cmd{Draft\_AddToLayer} \\
注释样式 & Annotation Styles & \cmd{Draft\_AnnotationStyleEditor} \\
应用当前样式 & Apply Current Style & \cmd{Draft\_ApplyStyle} \\
圆弧 & Arc & \cmd{Draft\_Arc} \\
三点创建圆弧 & Arc From 3 Points & \cmd{Draft\_Arc\_3Points} \\
阵列 & Array & \cmd{Draft\_OrthoArray} \\
自动分组 & Auto-Group & \cmd{Draft\_AutoGroup} \\
贝塞尔曲线 & B-Spline & \cmd{Draft\_BSpline} \\
贝塞尔曲线 & Bézier Curve & \cmd{Draft\_BezCurve} \\
圆 & Circle & \cmd{Draft\_Circle} \\
环形阵列 & Circular Array & \cmd{Draft\_CircularArray} \\
克隆 & Clone & \cmd{Draft\_Clone} \\
转换多段线/B样条曲线 & Convert Wire/B-Spline & \cmd{Draft\_WireToBSpline} \\
三次贝塞尔曲线 & Cubic Bézier Curve & \cmd{Draft\_CubicBezCurve} \\
尺寸 & Dimension & \cmd{Draft\_Dimension} \\
降级 & Downgrade & \cmd{Draft\_Downgrade} \\
草绘转草图 & Draft to Sketch & \cmd{Draft\_Draft2Sketch} \\
编辑 & Edit & \cmd{Draft\_Edit} \\
椭圆 & Ellipse & \cmd{Draft\_Ellipse} \\
面引用连接 & Facebinder & \cmd{Draft\_Facebinder} \\
圆角 & Fillet & \cmd{Draft\_Fillet} \\
翻转标注 & Flip Dimension & \cmd{Draft\_FlipDimension} \\
强制更新 2D 视图 & Force 2D View Update & \cmd{Draft\_UpdateShape2DView} \\
剖面线 & Hatch & \cmd{Draft\_Hatch} \\
\bottomrule
\end{longtable}
\section{工程图（TechDraw）}
\begin{longtable}{@{}p{0.32\linewidth}p{0.30\linewidth}p{0.30\linewidth}@{}}
\toprule 中文 & English & 命令 ID \\ \midrule \endhead
对齐水平链式尺寸 & Align Horizontal Chain Dimensions & \cmd{TechDraw\_ExtensionPosChainDimensionGroup} \\
级联水平尺寸 & Cascade Horizontal Dimensions & \cmd{TechDraw\_ExtensionCascadeDimensionGroup} \\
中心线 & Centerline & \cmd{TechDraw\_CenterLineGroup} \\
圆中心线 & Circle Centerlines & \cmd{TechDraw\_ExtensionCircleCenterLinesGroup} \\
装饰单点圆 & Cosmetic 1 Point Circle & \cmd{TechDraw\_ExtensionDrawCirclesGroup} \\
装饰性交点顶点 & Cosmetic Intersection Vertices & \cmd{TechDraw\_CommandVertexCreationGroup} \\
装饰平行线 & Cosmetic Parallel Line & \cmd{TechDraw\_ExtensionLinePPGroup} \\
装饰性螺纹孔侧视图 & Cosmetic Thread Hole Side View & \cmd{TechDraw\_ExtensionThreadsGroup} \\
修饰顶点 & Cosmetic Vertex & \cmd{TechDraw\_CosmeticVertexGroup} \\
尺寸 & Dimension & \cmd{TechDraw\_CompDimensionTools} \\
延长线 & Extend Line & \cmd{TechDraw\_ExtensionExtendShortenLineGroup} \\
扩展尺寸 & Extent Dimension & \cmd{TechDraw\_ExtentGroup} \\
水平链式标注 & Horizontal Chain Dimension & \cmd{TechDraw\_ExtensionCreateChainDimensionGroup} \\
水平倒角标注 & Horizontal Chamfer Dimension & \cmd{TechDraw\_ExtensionChamferDimensionGroup} \\
水平坐标标注 & Horizontal Coordinate Dimension & \cmd{TechDraw\_ExtensionCreateCoordDimensionGroup} \\
增加小数位数 & Increase Decimal Places & \cmd{TechDraw\_ExtensionIncreaseDecreaseGroup} \\
插入 '⌀' 前缀 & Insert '⌀' Prefix & \cmd{TechDraw\_ExtensionInsertPrefixGroup} \\
剖面视图（简单或复杂） & Section View (Simple or Complex) & \cmd{TechDraw\_SectionGroup} \\
视图堆叠顺序 & View Stacking Order & \cmd{TechDraw\_StackGroup} \\
活动视图 & Active View & \cmd{TechDraw\_ActiveView} \\
将视图添加到裁剪组 & Add View To Clip Group & \cmd{TechDraw\_ClipGroupAdd} \\
对齐倾斜链式尺寸 & Align Oblique Chain Dimensions & \cmd{TechDraw\_ExtensionPosObliqueChainDimension} \\
对齐垂直链式尺寸 & Align Vertical Chain Dimensions & \cmd{TechDraw\_ExtensionPosVertChainDimension} \\
水平对齐顶点/边 & Align Vertices/Edge Horizontally & \cmd{TechDraw\_AlignVertexesHorizontally} \\
垂直对齐顶点/边 & Align Vertices/Edge Vertically & \cmd{TechDraw\_AlignVertexesVertically} \\
角度尺寸 & Angle Dimension & \cmd{TechDraw\_AngleDimension} \\
基于3点的角度标注 & Angle Dimension From 3 Points & \cmd{TechDraw\_3PtAngleDimension} \\
弧长标注 & Arc Length Annotation & \cmd{TechDraw\_ExtensionArcLengthAnnotation} \\
弧长标注 & Arc Length Dimension & \cmd{TechDraw\_ExtensionCreateLengthArc} \\
面积注释 & Area Annotation & \cmd{TechDraw\_AreaDimension} \\
\bottomrule
\end{longtable}
\section{装配（Assembly）}
\begin{longtable}{@{}p{0.32\linewidth}p{0.30\linewidth}p{0.30\linewidth}@{}}
\toprule 中文 & English & 命令 ID \\ \midrule \endhead
齿轮/皮带配合 & Gears/Belt Joint & \cmd{Assembly\_CreateJointGearBelt} \\
插入零件 & Insert Component & \cmd{Assembly\_Insert} \\
激活装配体 & Activate Assembly & \cmd{Assembly\_ActivateAssembly} \\
角度配合 & Angle Joint & \cmd{Assembly\_CreateJointAngle} \\
球窝配合 & Ball Joint & \cmd{Assembly\_CreateJointBall} \\
皮带配合 & Belt Joint & \cmd{Assembly\_CreateJointBelt} \\
材料清单 & Bill of Materials & \cmd{Assembly\_CreateBom} \\
创建刚性组 & Create Rigid Group & \cmd{Assembly\_CreateJointRigidGroup} \\
同轴配合 & Cylindrical Joint & \cmd{Assembly\_CreateJointCylindrical} \\
距离配合 & Distance Joint & \cmd{Assembly\_CreateJointDistance} \\
爆炸视图 & Exploded View & \cmd{Assembly\_CreateView} \\
导出 ASMT 文件 & Export ASMT File & \cmd{Assembly\_ExportASMT} \\
固定配合 & Fixed Joint & \cmd{Assembly\_CreateJointFixed} \\
齿轮配合 & Gears Joint & \cmd{Assembly\_CreateJointGears} \\
转到关联的装配体 & Go to Linked Assembly & \cmd{Assembly\_LinkSelectLinked} \\
新建装配体 & New Assembly & \cmd{Assembly\_CreateAssembly} \\
新零件 & New Part & \cmd{Assembly\_InsertNewPart} \\
平行配合 & Parallel Joint & \cmd{Assembly\_CreateJointParallel} \\
垂直配合 & Perpendicular Joint & \cmd{Assembly\_CreateJointPerpendicular} \\
齿轮齿条配合 & Rack and Pinion Joint & \cmd{Assembly\_CreateJointRackPinion} \\
旋转配合 & Revolute Joint & \cmd{Assembly\_CreateJointRevolute} \\
螺纹配合 & Screw Joint & \cmd{Assembly\_CreateJointScrew} \\
选择零件配合 & Select Component Joints & \cmd{Assembly\_SelectJointsOfComponent} \\
选择具有自由度的零件 & Select Components With DoFs & \cmd{Assembly\_SelectComponentsWithDoFs} \\
选择冲突约束 & Select Conflicting Constraints & \cmd{Assembly\_SelectConflictingConstraints} \\
选择格式错误的约束 & Select Malformed Constraints & \cmd{Assembly\_SelectMalformedConstraints} \\
选择冗余约束 & Select Redundant Constraints & \cmd{Assembly\_SelectRedundantConstraints} \\
模拟 & Simulation & \cmd{Assembly\_CreateSimulation} \\
滑块配合 & Slider Joint & \cmd{Assembly\_CreateJointSlider} \\
快照 & Snapshot & \cmd{Assembly\_CreateSnapshot} \\
\bottomrule
\end{longtable}
\section{电子表格（Spreadsheet）}
\begin{longtable}{@{}p{0.32\linewidth}p{0.30\linewidth}p{0.30\linewidth}@{}}
\toprule 中文 & English & 命令 ID \\ \midrule \endhead
底部对齐 & Align Bottom & \cmd{Spreadsheet\_AlignBottom} \\
水平居中对齐 & Align Horizontal Center & \cmd{Spreadsheet\_AlignCenter} \\
左对齐 & Align Left & \cmd{Spreadsheet\_AlignLeft} \\
右对齐 & Align Right & \cmd{Spreadsheet\_AlignRight} \\
顶部对齐 & Align Top & \cmd{Spreadsheet\_AlignTop} \\
垂直居中对齐 & Align Vertical Center & \cmd{Spreadsheet\_AlignVCenter} \\
粗体文本 & Bold Text & \cmd{Spreadsheet\_StyleBold} \\
导出电子表格 & Export Spreadsheet & \cmd{Spreadsheet\_Export} \\
导入电子表格 & Import Spreadsheet & \cmd{Spreadsheet\_Import} \\
斜体文本 & Italic Text & \cmd{Spreadsheet\_StyleItalic} \\
合并单元格 & Merge Cells & \cmd{Spreadsheet\_MergeCells} \\
新建电子表格 & New Spreadsheet & \cmd{Spreadsheet\_CreateSheet} \\
设置别名 & Set Alias & \cmd{Spreadsheet\_SetAlias} \\
拆分单元格 & Split Cell & \cmd{Spreadsheet\_SplitCell} \\
下划线文本 & Underline Text & \cmd{Spreadsheet\_StyleUnderline} \\
\bottomrule
\end{longtable}
\section{加工（CAM）}
\begin{longtable}{@{}p{0.32\linewidth}p{0.30\linewidth}p{0.30\linewidth}@{}}
\toprule 中文 & English & 命令 ID \\ \midrule \endhead
修饰操作 & Dressup Operations & \cmd{CAM\_DressupTools} \\
钻孔操作 & Drilling Operations & \cmd{CAM\_DrillingTools} \\
雕刻加工 & Engraving Operations & \cmd{CAM\_EngraveTools} \\
后处理加工 & Post process Operations & \cmd{CAM\_PostTools} \\
模拟器 & Simulators & \cmd{CAM\_SimTools} \\
三维工具 & 3dTools & \cmd{CAM\_3dTools} \\
自适应 & Adaptive & \cmd{CAM\_Adaptive} \\
添加刀具… & Add Toolbit... & \cmd{CAM\_ToolBitDock} \\
区域 & Area & \cmd{CAM\_Area} \\
区域工作平面 & Area Workplane & \cmd{CAM\_Area\_Workplane} \\
阵列 & Array & \cmd{CAM\_Array} \\
轴映射 & Axis Map & \cmd{CAM\_DressupAxisMap} \\
边界 & Boundary & \cmd{CAM\_DressupPathBoundary} \\
边界2 & Boundary2 & \cmd{CAM\_DressupPathBoundary2} \\
CAM模拟器 & CAM Simulator & \cmd{CAM\_SimulatorGL} \\
CAMotics & CAMotics & \cmd{CAM\_Camotics} \\
注释 & Comment & \cmd{CAM\_Comment} \\
组合 & Compound & \cmd{CAM\_Compound} \\
复制 & Copy & \cmd{CAM\_Copy} \\
复制操作 & Copy Operation & \cmd{CAM\_OperationCopy} \\
自定义 & Custom & \cmd{CAM\_Custom} \\
去毛刺 & Deburr & \cmd{CAM\_Deburr} \\
狗骨圆角 & Dogbone & \cmd{CAM\_DressupDogbone} \\
拖刀 & Drag Knife & \cmd{CAM\_DressupDragKnife} \\
钻孔 & Drilling & \cmd{CAM\_Drilling} \\
雕刻 & Engrave & \cmd{CAM\_Engrave} \\
导出模板 & Export Template & \cmd{CAM\_ExportTemplate} \\
完成所选循环 & Finish Selecting Loop & \cmd{CAM\_SelectLoop} \\
夹具 & Fixture & \cmd{CAM\_Fixture} \\
从形状 & From Shape & \cmd{CAM\_Shape} \\
\bottomrule
\end{longtable}
\section{网格（Mesh）}
\begin{longtable}{@{}p{0.32\linewidth}p{0.30\linewidth}p{0.30\linewidth}@{}}
\toprule 中文 & English & 命令 ID \\ \midrule \endhead
添加三角形 & Add Triangle & \cmd{Mesh\_AddFacet} \\
边界框信息 & Bounding Box Info & \cmd{Mesh\_BoundingBox} \\
封闭孔 & Close Hole & \cmd{Mesh\_FillInteractiveHole} \\
横截面 & Cross-Sections & \cmd{Mesh\_CrossSections} \\
曲率信息 & Curvature Info & \cmd{Mesh\_CurvatureInfo} \\
曲率图 & Curvature Plot & \cmd{Mesh\_VertexCurvature} \\
剪切 & Cut & \cmd{Mesh\_PolyCut} \\
精简 & Decimate & \cmd{Mesh\_Decimating} \\
差集 & Difference & \cmd{Mesh\_Difference} \\
评估和修复 & Evaluate and Repair & \cmd{Mesh\_Evaluation} \\
评估实体 & Evaluate Solid & \cmd{Mesh\_EvaluateSolid} \\
导出网格… & Export Mesh... & \cmd{Mesh\_Export} \\
面信息 & Face Info & \cmd{Mesh\_EvaluateFacet} \\
填充孔 & Fill Holes & \cmd{Mesh\_FillupHoles} \\
翻转法线 & Flip Normals & \cmd{Mesh\_FlipNormals} \\
协调法线 & Harmonize Normals & \cmd{Mesh\_HarmonizeNormals} \\
导入网格… & Import Mesh... & \cmd{Mesh\_Import} \\
交集 & Intersection & \cmd{Mesh\_Intersection} \\
合并 & Merge & \cmd{Mesh\_Merge} \\
来自几何体的网格 & Mesh From Geometry & \cmd{Mesh\_FromGeometry} \\
来自形状的网格 & Mesh From Shape & \cmd{Mesh\_FromPartShape} \\
精制 & Refinement & \cmd{Mesh\_RemeshGmsh} \\
常规实体 & Regular Solid & \cmd{Mesh\_BuildRegularSolid} \\
删除组件 & Remove Components & \cmd{Mesh\_RemoveComponents} \\
手动移除组件 & Remove Components Manually & \cmd{Mesh\_RemoveCompByHand} \\
缩放 & Scale & \cmd{Mesh\_Scale} \\
平面剖切 & Section From Plane & \cmd{Mesh\_SectionByPlane} \\
段 & Segment & \cmd{Mesh\_PolySegm} \\
段 & Segmentation & \cmd{Mesh\_Segmentation} \\
从最佳拟合曲面进行分割 & Segmentation From Best-Fit Surfaces & \cmd{Mesh\_SegmentationBestFit} \\
\bottomrule
\end{longtable}
\section{BIM（BIM）}
\begin{longtable}{@{}p{0.32\linewidth}p{0.30\linewidth}p{0.30\linewidth}@{}}
\toprule 中文 & English & 命令 ID \\ \midrule \endhead
圆弧工具 & Arc Tools & \cmd{BIM\_ArcTools} \\
阵列工具 & Array Tools & \cmd{BIM\_ArrayTools} \\
轴线工具 & Axis Tools & \cmd{BIM\_AxisTools} \\
布尔工具 & Boolean Tools & \cmd{BIM\_BooleanTools} \\
克隆工具 & Cloning Tools & \cmd{BIM\_CloneTools} \\
创建 2D 视图 & Create 2D Views & \cmd{BIM\_Create2DViews} \\
拉伸工具 & Extrude Tools & \cmd{BIM\_ExtrudeTools} \\
通用 3D 工具 & Generic 3D Tools & \cmd{BIM\_GenericTools} \\
IFC 管理 & IFC Management & \cmd{BIM\_IfcManageTools} \\
偏移工具 & Offset Tools & \cmd{BIM\_OffsetTools} \\
报表工具 & Report Tools & \cmd{BIM\_ReportTools} \\
样条曲线工具 & Spline Tools & \cmd{BIM\_SplineTools} \\
2D 绘图 & 2D Drawing & \cmd{BIM\_DrawingView} \\
2D 偏移 & 2D Offset & \cmd{BIM\_Offset2D} \\
对齐标注 & Aligned Dimension & \cmd{BIM\_DimensionAligned} \\
梁 & Beam & \cmd{BIM\_Beam} \\
BIM 示例 & BIM Examples & \cmd{BIM\_Examples} \\
BIM 帮助 & BIM Help & \cmd{BIM\_Help} \\
BIM 设置 & BIM Setup & \cmd{BIM\_Setup} \\
BIM 教程 & BIM Tutorial & \cmd{BIM\_Tutorial} \\
BIM欢迎屏幕 & BIM Welcome Screen & \cmd{BIM\_Welcome} \\
立方体 & Box & \cmd{BIM\_Box} \\
克隆 & Clone & \cmd{BIM\_Clone} \\
柱 & Column & \cmd{BIM\_Column} \\
组合 & Compound & \cmd{BIM\_Compound} \\
转换为 BIM & Convert to BIM & \cmd{BIM\_Convert} \\
复制 & Copy & \cmd{BIM\_Copy} \\
覆盖层 & Covering & \cmd{BIM\_Covering} \\
差集 & Difference & \cmd{BIM\_Cut} \\
门 & Door & \cmd{BIM\_Door} \\
\bottomrule
\end{longtable}
\section{曲面（Surface）}
\begin{longtable}{@{}p{0.32\linewidth}p{0.30\linewidth}p{0.30\linewidth}@{}}
\toprule 中文 & English & 命令 ID \\ \midrule \endhead
混合曲线 & Blend Curve & \cmd{Surface\_BlendCurve} \\
网格上的曲线 & Curve on Mesh & \cmd{Surface\_CurveOnMesh} \\
延伸面 & Extend Face & \cmd{Surface\_ExtendFace} \\
填充边界曲线 & Fill Boundary Curves & \cmd{Surface\_GeomFillSurface} \\
填充 & Filling & \cmd{Surface\_Filling} \\
截面 & Sections & \cmd{Surface\_Sections} \\
曲面切割 & Surface Cut & \cmd{Surface\_Cut} \\
\bottomrule
\end{longtable}

\chapter{中英术语对照}
\markboth{附录 B 中英术语对照}{附录 B 中英术语对照}

MobileCAD 的中文界面直接采用 FreeCAD 官方译法。对照本表，可以无缝切换到英文界面或桌面 FreeCAD。

\section{建模概念}

\begin{longtable}{@{}p{0.16\linewidth}p{0.24\linewidth}p{0.52\linewidth}@{}}
\toprule \textbf{中文} & \textbf{English} & \textbf{说明} \\ \midrule \endhead
参数化建模 & parametric modelling & 用尺寸、约束与特征历史描述模型，改参数即改模型 \\
草图 & sketch & 2D 轮廓，通常作为 3D 特征的截面 \\
约束 & constraint & 限定几何关系的规则（水平、相切、距离……） \\
自由度 & degree of freedom (DoF) & 草图未被约束的独立变量数 \\
完全约束 & fully constrained & 自由度为 0，形状被唯一确定 \\
构造几何 & construction geometry & 只辅助作图、不参与成形的线 \\
尺寸 & dimension & 数值约束（长度、角度、半径……） \\
参考（从动） & reference / driven & 只测量不驱动的尺寸 \\
凸台 & Pad & 把闭合轮廓沿法向拉伸出实体 \\
凹槽 & Pocket & 用轮廓在实体上挖出凹坑/通孔 \\
旋转 & Revolution & 轮廓绕轴旋转成回转体 \\
槽 & Groove & 在回转体上挖环形槽 \\
放样 & Loft & 用多个截面过渡成实体 \\
管道 / 扫掠 & Pipe / Sweep & 轮廓沿路径扫掠 \\
螺旋 & Helix & 沿螺旋线扫掠（弹簧、螺纹） \\
圆角 & Fillet & 把尖边变成圆弧过渡 \\
倒角 & Chamfer & 把尖边切出斜面 \\
拔模 & Draft & 给侧面加脱模斜度 \\
厚度 & Thickness & 把实体抽壳成薄壁 \\
孔 & Hole & 标准孔特征（可带螺纹、沉头） \\
阵列 & Pattern & 线性/极坐标/路径/点阵列 \\
镜像 & Mirrored & 关于平面镜像 \\
实体（体） & Body & PartDesign 里承载特征历史的容器 \\
特征 & feature & 建模历史中的一步（凸台、圆角……）\\
Tip & Tip & 实体“计算到哪一步”的标记 \\
基准面/轴/点 & Datum plane/axis/point & 辅助建模的参考几何 \\
附着 & Attachment & 把草图/基准几何挂到面、边或坐标上 \\
布尔运算 & boolean operation & 并、差、交 \\
B-rep & boundary representation & 用边界曲面描述实体的几何表示 \\
\bottomrule
\end{longtable}

\section{装配、工程图与仿真}

\begin{longtable}{@{}p{0.16\linewidth}p{0.24\linewidth}p{0.52\linewidth}@{}}
\toprule \textbf{中文} & \textbf{English} & \textbf{说明} \\ \midrule \endhead
装配 & Assembly & 多个构件的组合与约束 \\
构件 & component & 装配中的零件/子装配 \\
关节 / 配合 & joint & 限定两个构件相对运动的关系 \\
接地 & grounded & 把构件固定在空间 \\
求解 & solve & 计算装配的自由度与位置 \\
工程图 & drawing (TechDraw) & 由 3D 生成的二维图纸 \\
投影组 & projection group & 前视/上视/右视等一组视图 \\
剖视图 & section view & 用剖切面切开显示的视图 \\
局部详图 & detail view & 放大局部区域 \\
尺寸标注 & dimension & 图纸上的尺寸 \\
电子表格 & spreadsheet & 参数表，可被表达式引用 \\
别名 & alias & 单元格的名字，如 \file{Thickness} \\
表达式 & expression & 以 \file{=} 开头的公式 \\
单位制 & unit system & 毫米/米、英寸等 \\
有限元 & FEM & 有限元分析 \\
网格划分 & meshing & 把实体离散为单元 \\
边界条件 & boundary condition & 固定、载荷等 \\
求解器 & solver & CalculiX（ccx）等 \\
后处理 & post-processing & 云图、变形、导出结果 \\
刀路 & toolpath & CAM 生成的加工路径 \\
后处理器 & post-processor & 把刀路转成机床 G-code \\
\bottomrule
\end{longtable}

\section{文件与技术}

\begin{longtable}{@{}p{0.16\linewidth}p{0.24\linewidth}p{0.52\linewidth}@{}}
\toprule \textbf{中文} & \textbf{English} & \textbf{说明} \\ \midrule \endhead
FCStd & FreeCAD document & 原生文档，保留建模历史 \\
STEP & STEP (.step/.stp) & 通用的实体交换格式 \\
IGES & IGES (.iges/.igs) & 较老的曲面/线框交换格式 \\
STL & STL (.stl) & 三角网格，3D 打印常用 \\
OBJ & OBJ (.obj) & 网格格式，带材质/贴图 \\
3MF & 3MF (.3mf) & 现代 3D 打印格式 \\
DXF / DWG & DXF / DWG & 二维 CAD 交换格式 \\
SVG & SVG & 矢量图（工程图/导入图形）\\
PDF & PDF & 工程图输出/打印 \\
IFC & IFC & 建筑信息模型交换格式 \\
G-code & G-code & 数控机床指令 \\
SAF & Storage Access Framework & Android 系统文件选择器 \\
手势 & gesture & 触屏操作 \\
精密指针 & precision pointer & 长按时出现在手指上方的准星 \\
导航立方体 & NaviCube & 右上角的方向指示与视图切换器 \\
事务 & transaction & 一次可撤销的修改单位 \\
恢复副本 & recovery copy & 崩溃后自动保存的文档副本 \\
\bottomrule
\end{longtable}

\chapter{常见问题}
\markboth{附录 C 常见问题}{附录 C 常见问题}

\begin{description}[style=nextline,leftmargin=0pt]

\item[它和桌面 FreeCAD 是什么关系？]
同一个引擎、同一种文档格式。MobileCAD 重写了界面（Kotlin + Compose），并用 Python 模拟 Qt，让 FreeCAD
自带的任务面板能在触屏上运行。几何算法、文件格式与 FreeCAD 完全一致。

\item[能替代电脑上的 CAD 吗？]
对“设计—修改—出图”的中等复杂度工作，可以。超大装配（上千个零件）、复杂曲面重建、大规模有限元，
桌面版仍更合适。建议的分工：手机/平板做构思、修改、现场测量与审图；桌面做重计算与批量处理。

\item[支持 DWG 吗？]
不支持。DWG 需要外部转换器（ODA/LibreDWG），未打包进 APK。请把 DWG 另存为 \textbf{DXF} 再导入。

\item[有 Python 控制台和插件吗？]
没有交互式控制台，也没有插件管理器。但\textbf{宏仍可运行}（\cmd{Std\_DlgMacroExecute}），
FreeCAD 的 Python 命令也照常工作——BIM/CAM/FEM 的界面本身就是 Python 写的。

\item[文件存在哪里？会自动保存吗？]
位置由你通过系统文件选择器决定，应用不强制目录。自动保存是“崩溃恢复副本”，不是正式保存：
每 60 秒写一份到应用私有目录，异常退出后自动恢复。请仍然养成手动保存的习惯。

\item[能同时打开两个文档吗？]
不能。MobileCAD 是单文档界面，新开文件会替换当前视图。

\item[为什么有些菜单项是灰的？]
那些是桌面窗口管理类命令（多窗口、停靠面板、工具栏定制、宏录制等），移动端没有对应物。点按会提示
“该命令尚未在移动端提供”。

\item[中文界面怎么开启？]
默认跟随系统：系统是中文就显示简体中文。也可在 \menupath{编辑 → 偏好设置 → 语言} 里固定为
\ui{简体中文}、\ui{英语} 或 \ui{系统语言}。

\item[支持哪些设备？]
Android 8.0 及以上、arm64 架构。手机、平板均可；Quest 等头显使用 MobileCAD XR 版本。安装包约 165 MB，
运行时约 1 GB。

\item[能离线使用吗？]
能。建模、装配、出图、有限元、CAM 全部本地完成。只有“检查更新”、打开网页帮助/捐赠链接时才联网。

\item[有限元能算多大？]
取决于设备内存与耐心。手机上建议单元数控制在几十万以内；先用粗网格验证，再局部加密。
只支持 CalculiX 求解器，Netgen 网格器不可用。

\item[能直接打印图纸吗？]
工程图可导出 PDF/SVG/DXF；移动端没有系统打印对话框，“打印”等同于导出当前页 PDF。

\item[GPL 许可证可以商用吗？]
可以。GPL-3.0 允许商业使用，但如果你分发修改后的程序，需要同时提供对应源代码。内部使用不受限制。

\item[怎样反馈问题或贡献翻译？]
在 \url{https://github.com/LucasCoraca/mobilecad} 提 issue（附设备型号、Android 版本、应用版本、
FreeCAD 提交号、复现步骤与截图）。翻译词表在 \file{app/app/src/main/assets/i18n/}，
中文条目可以直接提交 PR。

\end{description}


\end{document}
